\documentclass[12pt,a4paper]{article}

\usepackage[utf8]{inputenc}
\usepackage[T1]{fontenc}
\usepackage{lmodern}
\usepackage{geometry}
\usepackage{setspace}
\usepackage{booktabs}
\usepackage{longtable}
\usepackage{array}
\usepackage{enumitem}
\usepackage{amssymb}
\usepackage{xcolor}
\usepackage{hyperref}

\geometry{margin=1in}
\setlength{\emergencystretch}{3em}
% Permite quebrar caminhos e nomes de variáveis em / _ . -
\makeatletter\g@addto@macro\UrlBreaks{\do\_\do\-\do\.}\makeatother
\onehalfspacing
\hypersetup{
    colorlinks=true,
    linkcolor=blue,
    urlcolor=blue
}

% Status dos resultados no guia rápido
\newcommand{\robusto}{\textcolor{green!45!black}{\textbf{robusto}}}
\newcommand{\fragil}{\textcolor{orange!85!black}{\textbf{frágil}}}
\newcommand{\superado}{\textcolor{red!70!black}{\textbf{superado}}}
\newcommand{\aberto}{\textcolor{blue!70!black}{\textbf{em aberto}}}

\title{\textbf{Documentação do Projeto FTPL}\\
\large Denominação da Dívida Pública Brasileira, Respaldo Fiscal e Inflação}
\author{Documento de acompanhamento do pipeline empírico}
\date{\today}

\begin{document}

\maketitle

\begin{abstract}
Este documento registra o estado atual do projeto empírico sobre denominação da
dívida pública brasileira, teoria fiscal do nível de preços (FTPL) e inflação. A
Parte~I é um guia rápido para quem chega ao projeto: a pergunta, a cadeia empírica,
o que já sabemos, as decisões pendentes e os próximos passos. A Parte~II descreve a
estratégia e os resultados na ordem lógica: modelo, equação de previsão das
emissões, purga da notícia fiscal, projeções locais, o problema de identificação
dentro do mês e as surpresas de oferta usadas como instrumento. A Parte~III é a
referência de dados: a cadeia de scripts e painéis, as fontes, as convenções de
nomes e um glossário. Os apêndices guardam detalhes de construção, auditorias e o
histórico de correções.
\end{abstract}

\newpage

\tableofcontents

\newpage

%======================================================================
\part{Guia rápido}
%======================================================================

\section{Onde estamos}\label{sec:guia}

\subsection{A pergunta em uma página}

O projeto investiga se a denominação da dívida pública brasileira altera a
transmissão de choques de emissão para a inflação e para preços de ativos. Pela
FTPL, o valor real das obrigações do governo é igual ao valor presente esperado dos
superávits primários. Se uma emissão não é respaldada por superávits futuros, algo
precisa se ajustar para reduzir o valor real das obrigações. O nível de preços só
consegue reduzir o valor real da dívida cujo valor depende dele:

\begin{itemize}[leftmargin=1.2cm]
    \item \(N\), prefixados (LTN e NTN-F): o nível de preços corrói todo o valor real
    (elasticidade \(\kappa^N=1\));
    \item \(F\), flutuantes (LFT): corrosão parcial, que depende de quanto a Selic
    reage à inflação (\(0<\kappa^F<1\));
    \item \(R\), indexados à inflação (NTN-B): nenhuma corrosão (\(\kappa^R=0\)).
\end{itemize}

Daí a previsão testável: para emissões inesperadas do mesmo tamanho e não
explicadas por notícia fiscal, a resposta da inflação deve ser ordenada
\(|\beta^N|>|\beta^F|>|\beta^R|\). A comparação \(N\)-\(R\) é o teste principal,
porque as duas elasticidades são fixadas pelo contrato; a ordenação completa
depende do regime monetário. Uma explicação baseada apenas em restrição de recursos
prevê respostas parecidas entre denominações; a diferença entre elas é o conteúdo
específico do teste FTPL (Seção~\ref{sec:teoria}).

\subsection{A cadeia empírica}

\begin{center}
\small
\begin{tabular}{p{0.04\textwidth}p{0.40\textwidth}p{0.44\textwidth}}
\toprule
 & Etapa & Objeto e script \\
\midrule
1 & Emissão bruta mensal por classe, sobre o PIB de 12 meses defasado &
\(g_t^j\), \path{g_{j}_pib}; notebook de dados \\
\(\downarrow\) & & \\
2 & Equação de previsão com informação de \(t-1\) (vencimentos programados, curva,
macro, PAF, Focus) & inovação \(u_t^j\), \path{u_{j}_base};
\path{build_issuance_forecasts.py} (Seção~\ref{sec:prev}) \\
\(\downarrow\) & & \\
3 & Purga da notícia fiscal do mês (revisões Focus do primário) &
choque \(\varepsilon_t^j\), \path{eps_{j}_base_fprim}; \path{build_purged_shocks.py}
(Seção~\ref{sec:purga}) \\
\(\downarrow\) & & \\
4 & Projeções locais dos resultados nos três choques juntos & \(\beta_h^j\),
\(h=0,\dots,24\); \path{build_local_projections.py} (Seção~\ref{sec:lp}) \\
\(\uparrow\) & & \\
5 & Instrumento: surpresas no lote ofertado nos leilões, com informação anterior ao
edital & \path{offsurp_{j}_naive}; \path{build_offer_surprises.py}
(Seção~\ref{sec:prop1}) \\
\bottomrule
\end{tabular}
\end{center}

A etapa~5 existe porque a etapa~2 só usa informação de \(t-1\): ela garante que
\(u_t^j\) não seja previsível com o passado, mas não que seja independente do
mercado dentro do mês \(t\). O Tesouro emite mais quando o mercado do mês está
favorável, o que vicia as projeções por MQO (Seção~\ref{sec:ident}).

\subsection{O que sabemos hoje}

\begin{center}
\small
\begin{longtable}{p{0.50\textwidth}p{0.14\textwidth}p{0.24\textwidth}}
\toprule
Resultado & Status & Onde \\
\midrule
\endhead
A equação de previsão explica 59--71\% da emissão; vencimentos de \(t\) elevam a
emissão da mesma classe (\(\approx 0{,}09\) por real) & \robusto &
Seção~\ref{sec:prev} \\
Revisões de pesquisa (Focus, Prisma) explicam quase nada das inovações; a escolha
de \(\Delta\mathcal{F}_t\) quase não altera os choques & \robusto &
Seção~\ref{sec:purga} \\
O choque de risco fiscal de Carvalho Filho (2026) explica cerca de 10\% de \(u^R\):
em meses de estresse fiscal o Tesouro emite menos NTN-B & \robusto, mas usa preços
de \(t\) & Seção~\ref{sec:cf} \\
Por MQO, a ordenação prevista falha: \(\beta^N-\beta^R<0\) no nível de preços &
\superado: resultado viesado & Seção~\ref{sec:lp_res} \\
Os choques \(N\) e \(R\) vêm com juros em queda, real apreciado e EMBI em queda no
próprio mês: emissão reage ao mercado do mês & \robusto & Seção~\ref{sec:ident} \\
Contados por lote, 78\% (\(N\)), 65\% (\(F\)) e 63\% (\(R\)) dos leilões vendem o
lote inteiro (a contagem por título, 28\% em \(R\), estava errada) & \robusto &
Seção~\ref{sec:ident}, Apêndice~\ref{app:hist} \\
Com as surpresas de oferta como instrumento, o viés em \(N\) desaparece;
\(\beta^N-\beta^R\) passa a \(+0{,}16\) (0,40) & \fragil: \(F\) de 1º estágio
\(\approx 9\) & Seção~\ref{sec:prop1} \\
Uma emissão inesperada de LFT de 0,1\% do PIB eleva o nível de preços em
\(\approx 0{,}2\%\) em 6--12 meses & \robusto em todas as versões IV e na forma
reduzida & Seção~\ref{sec:prop1} \\
O choque \(R\) continua sem identificação: instrumento fraco (\(F=6\)) e
ainda correlacionado com o mercado do mês & \aberto & Seção~\ref{sec:prop1} \\
A ordenação \(N>F>R\) não é rejeitada nem confirmada & \aberto & Seção~\ref{sec:prop1} \\
\bottomrule
\end{longtable}
\end{center}

\subsection{Decisões pendentes com o coautor}\label{sec:decisoes}

\begin{enumerate}[leftmargin=1.2cm]
    \item \textbf{Medida de notícia fiscal.} Confirmar o Focus primário de 12 e 36
    meses como \(\Delta\mathcal{F}_t\) principal, com Prisma, DLSP e Carvalho Filho
    como robustez (Seção~\ref{sec:purga_leitura}).
    \item \textbf{Purgar também a notícia fiscal precificada?} O choque de Carvalho
    Filho tira a reação do Tesouro ao estresse fiscal, mas traz preços de \(t\) para
    dentro dos choques.
    \item \textbf{Notícia fiscal como choque.} O relatório do coautor propõe usar a
    notícia fiscal como choque das projeções locais, em vez das inovações de emissão
    por classe. Isso muda a estratégia de identificação.
    \item \textbf{Identificação de \(R\).} Escolher entre buscar dados para a
    Proposta~2 (anúncios, alta frequência), refinar a Proposta~1 no nível do lote ou
    aceitar \(R\) como não identificado e concentrar o teste em \(N\)-\(F\)
    (Seção~\ref{sec:prop1_pend}).
    \item \textbf{Tratamento de 2020} (dummies na purga ou na projeção local).
\end{enumerate}

\subsection{Próximos passos}\label{sec:proximos}

\begin{enumerate}[leftmargin=1.2cm]
    \item Inferência robusta a instrumento fraco (Anderson--Rubin por subvetor) para
    as diferenças entre classes.
    \item Levantar as datas reais dos editais de leilão (hoje supostos na véspera) e,
    se possível, o volume total de propostas por leilão.
    \item Uma fonte de variação de oferta para \(R\) que não acompanhe a demanda.
    \item Proposta~3: limites para \(\beta^N-\beta^R\) sob valores plausíveis do viés.
    \item Interações com a participação defasada e dependência de estado pela
    ociosidade; bootstrap em blocos que repete a construção dos choques.
    \item Dados prometidos ou a obter: extensão da ETTJ (antes de 2006 e depois de
    jan/2025), planilhas originais do Prisma e exportação Qlik de 2024 em diante,
    datas de coleta do Prisma e de divulgação do primário.
    \item Estender o cronograma mensal de vencimentos para antes de 2006-11.
\end{enumerate}

\subsection{Como ler este documento}

\begin{itemize}[leftmargin=1.2cm]
    \item Para entender a estratégia e os resultados: Parte~II, na ordem.
    \item Para rodar o código ou localizar uma variável: Seção~\ref{sec:pipeline}
    (cadeia de scripts e comandos) e Seção~\ref{sec:convencoes} (nomes e glossário).
    \item Para a origem de cada bloco de dados: Seção~\ref{sec:dados}.
    \item Para detalhes históricos (o que foi corrigido e por quê): apêndices.
\end{itemize}

Documentos relacionados em \path{docs/}: \path{ModeloTeorico.tex} (modelo),
\path{FTPL_test.tex} (texto principal),
\path{ftpl_issuance_shocks_empirical_strategy_with_maturities.pdf} (nota
metodológica; as referências a ``eq.~\(n\) da nota'' apontam para ela),
\path{Nota_EmissoesBrutas.tex} (emissão bruta versus líquida), \path{Project.md}
(resumo em inglês da cadeia) e \path{ssrn-7702078.pdf} (Carvalho Filho, 2026).

%======================================================================
\part{Estratégia e resultados}
%======================================================================

\section{Modelo teórico e hipóteses testáveis}\label{sec:teoria}

A sequência é: a denominação da dívida determina elasticidades distintas de
reavaliação pelo nível de preços; essas elasticidades implicam uma ordenação das
respostas inflacionárias; e essa ordenação é testada com choques de emissão
construídos com informação predeterminada, purgados de notícia fiscal observável,
em projeções locais. O modelo completo está em \path{docs/ModeloTeorico.tex}.

\subsection{Valuation fiscal e denominação da dívida}

Considere três classes de passivos. Se \(\mathcal{V}_t^j\) é o valor real de mercado da
classe \(j\) e \(S_{t+\tau}\) é o superávit primário real, a restrição fiscal é
\begin{equation}
\mathcal{V}_t^N+\mathcal{V}_t^F+\mathcal{V}_t^R
=\mathbb{E}_t\sum_{\tau=0}^{\infty}m_{t,t+\tau}S_{t+\tau}
\equiv PV_t(S).
\label{eq:doc_valuation}
\end{equation}
A elasticidade de valuation em relação ao nível de preços é
\begin{equation}
\kappa_t^j\equiv-\frac{\partial\log\mathcal{V}_t^j}{\partial\log P_t},
\qquad
\kappa_t^N=1,\quad 0<\kappa_t^F<1,\quad \kappa_t^R=0.
\label{eq:doc_kappa}
\end{equation}
Os casos \(N\) e \(R\) são fixados pelo contrato; o caso \(F\) depende da reação da
Selic à inflação.

Considere um choque de emissão da classe \(j\). Se \(\rho_t^j\) é a parcela do
choque ainda respaldada por revisões residuais dos superávits esperados, o canal de
valuation implica
\begin{equation}
IRF_\pi^j(h)=\kappa_t^j(1-\rho_t^j)\Theta_\pi(h),
\qquad \Theta_\pi(h)>0.
\label{eq:doc_irf_theory}
\end{equation}
Com \(\rho_t^N=\rho_t^F=\rho_t^R<1\), segue
\(|IRF_\pi^N(h)|>|IRF_\pi^F(h)|>|IRF_\pi^R(h)|\). Com respaldo heterogêneo, a
condição relevante é \(\phi_t^N>\phi_t^F>\phi_t^R\), com
\(\phi_t^j\equiv\kappa_t^j(1-\rho_t^j)\). Por isso a comparação \(N\)-\(R\) é o
teste primário e a ordenação tripla é o teste mais forte e mais dependente do
regime. Uma explicação baseada somente em restrição de recursos, depois de
controlar impulso fiscal e ociosidade, prevê respostas aproximadamente semelhantes
entre denominações; a diferença entre elas é o conteúdo específico do teste FTPL.

\subsection{Da teoria à estimação}

\paragraph{Por que emissão bruta.} A medida de emissão é a emissão bruta no mercado
primário, normalizada pelo PIB acumulado em 12 meses defasado. A variação do
estoque não serve como choque, pois incorpora mecanicamente a capitalização da
Selic nas LFT e a correção pelo IPCA nas NTN-B, o que contaminaria os resultados de
inflação e juros (ver \path{docs/Nota_EmissoesBrutas.tex}).

\paragraph{Passo 1: inovação de emissão.} Separa-se a decisão de colocação de
refinanciamentos previsíveis e de condições conhecidas em \(t-1\):
\begin{equation}
g_t^j=\alpha^j+\delta_0^j m_{t|t-1}^j+\delta_1^j m_{t+1|t-1}^j
+\Lambda^{j\prime}Z_{t-1}+\mu_{m(t)}^j+u_t^j,
\label{eq:doc_forecast}
\end{equation}
em que \(m_{t|t-1}^j\) e \(m_{t+1|t-1}^j\) são o serviço da dívida programado para
\(t\) e \(t+1\), conhecido em \(t-1\), \(Z_{t-1}\) é o conjunto de informação
predeterminada e \(\mu_{m(t)}\) são dummies de mês (eq.~19 da nota;
Seção~\ref{sec:prev}).

\paragraph{Passo 2: purga da notícia fiscal.} Retira-se a parte da inovação
associada à notícia fiscal observável do mês:
\begin{equation}
u_t^j=\gamma^j\Delta\mathcal{F}_t+\eta^{j\prime}W_t+\varepsilon_t^j,
\qquad j\in\{N,F,R\},
\label{eq:doc_survey_purge}
\end{equation}
em que \(\Delta\mathcal{F}_t\) é a revisão da expectativa de superávit (eq.~20 da
nota) e \(W_t\) contém controles adicionais. O resíduo \(\varepsilon_t^j\) é lido,
de forma conservadora, como uma inovação não racionalizada por respaldo fiscal
observável, e não como choque perfeitamente desrespaldado (Seção~\ref{sec:purga}).

\paragraph{Passo 3: projeções locais.} Para um resultado \(y\) e horizonte \(h\),
\begin{equation}
y_{t+h}-y_{t-1}=\alpha_h+\beta_h^N\varepsilon_t^N
+\beta_h^F\varepsilon_t^F+\beta_h^R\varepsilon_t^R
+\Gamma_h'X_{t-1}+v_{t+h}
\label{eq:doc_local_projection}
\end{equation}
(eq.~30 da nota). \(\beta_h^{j,\pi}\) é a contrapartida empírica de
\(IRF_\pi^j(h)\), e as hipóteses são
\begin{equation}
H^{NR}: |\beta_h^{N,\pi}|>|\beta_h^{R,\pi}|,
\qquad
H^{NFR}: |\beta_h^{N,\pi}|>|\beta_h^{F,\pi}|>|\beta_h^{R,\pi}|.
\label{eq:doc_empirical_tests}
\end{equation}

\paragraph{Extensões previstas na nota.} Interações com a participação defasada
(eq.~36), dependência de estado pela ociosidade (eq.~37) e validação narrativa. O
instrumento de demanda regulatória por NTN-B, também previsto na nota, foi
descartado (Seção~\ref{sec:ident}).

\section{Equações de previsão de emissão}\label{sec:prev}

O script \path{src/build_issuance_forecasts.py} estima a eq.~(\ref{eq:doc_forecast}) para cada classe
\(j\in\{N,F,R\}\) e grava
\path{data/clean/ftpl_monthly_panel_with_issuance_innovations.csv}. A variável dependente é a emissão bruta sobre o PIB acumulado em 12 meses
defasado (\path{g_{j}_pib}). Os regressores são todos conhecidos antes dos leilões do
mês \(t\):

\begin{itemize}[leftmargin=1.2cm]
    \item \(m_{t|t-1}^j\): serviço da dívida (principal e cupons) que vence em \(t\)
    (\path{maturity_{j}_sched_pib_lag1}, Subseção~\ref{sec:venc});
    \item \(m_{t+1|t-1}^j\): serviço que vence em \(t+1\), calculado neste script a
    partir do estoque por título do RMD 2.2 em \(t-1\), com cupons indexados pelo IPCA
    ou IGP-M de \(t-2\), o último divulgado ao fim de \(t-1\). Antes de 2010-02 não há
    estoque por título em \(t-1\) e usa-se o cronograma de \(t+1\)
    (\path{maturity_lead1_source}). A correlação com o cronograma de \(t+1\) é de pelo
    menos 0,998 em \(N\), \(F\) e \(R\);
    \item defasagens 1 a 3 da emissão das três classes; estoque e participação da
    própria classe; janela de vencimentos de 12 meses; trocas, recompras e vendas e
    compras extraordinárias defasadas;
    \item Selic (defasagem 1), IPCA em 12 meses (defasagem 1, divulgado por volta do
    dia 10 de \(t\)), IBC-Br em 12 meses (defasagem 2, divulgado cerca de 45 dias
    depois) e variação do câmbio (defasagem 1);
    \item nível, inclinação e curvatura das curvas prefixada e IPCA no fim de \(t-1\);
    EMBI e volatilidade realizada da taxa prefixada de 5 anos em \(t-1\);
    \item expectativas Focus de primário e DLSP de 12 meses no fim de \(t-1\);
    \item gap do PAF da própria classe e gap da parcela vincenda em 12 meses (versão
    \path{rmd});
    \item dummies de mês do calendário.
\end{itemize}

As três equações usam a mesma amostra, necessária para a projeção local conjunta, e
os erros-padrão são HAC (Newey--West). A amostra base vai de 2006-11 a 2025-02 (220
meses, 46 parâmetros), limitada pelo cronograma mensal e pela ETTJ. O \(R^2\) é 0,71
em \(N\), 0,59 em \(F\) e 0,67 em \(R\), e o teste de Ljung--Box(12) não rejeita
ausência de autocorrelação (\(p\ge 0{,}36\)). Cada real que vence em \(t\) eleva a
emissão da mesma classe em cerca de 0,09 real, com significância em \(N\) e \(R\) e
não em \(F\); o vencimento de \(t+1\) não é significativo. A correlação entre as
inovações é 0,24 entre \(N\) e \(R\) e praticamente nula nos demais pares.

O resíduo \(u_t^j\) (\path{u_{j}_base}) é a inovação de emissão; é contemporâneo e
nunca entra como controle. As variantes de robustez são:

\begin{center}
\begin{tabular}{p{0.18\textwidth}p{0.72\textwidth}}
\toprule
Variante & Mudança em relação à base \\
\midrule
\path{long} & sem curva de juros; vai até 2026-06 (236 meses) \\
\path{w12} & janelas de 12 meses por classe no lugar do cronograma mensal; 2004-03 a 2026-06 (268 meses) \\
\path{pca} & curva resumida por componentes principais \\
\path{principal} & vencimentos apenas com principal; desde 2010-02 \\
\path{nolead} & sem \(m_{t+1|t-1}^j\) \\
\path{stock} & emissão e vencimentos sobre o estoque defasado da classe \\
\path{liq} & inclui a reserva de liquidez; desde 2015-02 \\
\path{ramp} & bloco real com NTN-B e NTN-C (somente a equação de \(R\)) \\
\path{base_rt} & erros de previsão em pseudo-tempo real, janela expansiva, desde 2014-07 \\
\bottomrule
\end{tabular}
\end{center}

A correlação das inovações da base com as variantes vai de 0,84 a 1,00, exceto em
\path{base_rt} (0,39 em \(N\), 0,71 em \(F\) e 0,72 em \(R\)), em que estimar 46
parâmetros apenas com dados até \(t-1\) gera muito ruído. Os diagnósticos estão em
\path{data/diagnostics/issuance_forecasts_*.csv} e a figura em
\path{data/figures/issuance_forecasts.png}.

\section{Escolha de \texorpdfstring{\(\Delta\mathcal{F}_t\)}{Delta F} e purga fiscal}\label{sec:purga}

A purga da eq.~(\ref{eq:doc_survey_purge}) exige escolher a medida de notícia fiscal
\(\Delta\mathcal{F}_t\). Esta seção registra as candidatas disponíveis no painel, os
prós e contras de cada uma, a comparação com a estratégia de Carvalho Filho (2026) e
os resultados da purga com a medida principal e as alternativas.

\subsection{Candidatas}

Todas as candidatas são notícia do mês \(t\) e entram apenas na purga. Nas revisões
de primário, o sinal positivo indica superávit esperado maior (notícia boa); nas
revisões de DLSP e no choque de Carvalho Filho, o sinal positivo é adverso. A tabela
resume a amostra base das inovações; as correlações com o mercado usam a variação no
mesmo mês do EMBI e do câmbio (alta = real mais fraco).

\begin{center}
\footnotesize
\begin{tabular}{p{0.27\textwidth}ccccc}
\toprule
Medida & Cobertura & Base & Zeros & AR(1) & \(\rho\) \(\Delta\)EMBI / câmbio \\
\midrule
Focus primário 12m & 2000-03 a 2026-11 & 220 & 25\% & 0,30 & \(-0{,}04\) / \(-0{,}11\) \\
Focus primário 36m & 2001-12 a 2026-11 & 220 & 30\% & 0,11 & 0,00 / 0,01 \\
Focus DLSP 12m & 2000-03 a 2026-11 & 220 & 4\% & 0,37 & 0,24 / 0,24 \\
Focus DLSP 36m & 2001-12 a 2026-11 & 220 & 5\% & 0,19 & 0,09 / 0,15 \\
Prisma primário 12m & 2016-01 a 2026-10 & 110 & 0\% & 0,34 & \(-0{,}15\) / \(-0{,}18\) \\
Prisma primário 36m & 2016-03 a 2026-10 & 99 & 0\% & 0,29 & 0,07 / \(-0{,}01\) \\
Surpresa do primário (Prisma) & 2016-01 a 2026-09 & 110 & 0\% & 0,03 & \(-0{,}03\) / \(-0{,}07\) \\
Choque de risco fiscal (Carvalho Filho) & 2005-12 a 2026-08 & 220 & -- & -- & -- \\
\bottomrule
\end{tabular}
\end{center}

O Focus primário de 12 meses e o Prisma primário de 12 meses têm correlação de 0,95.
As revisões de primário e de DLSP têm correlação entre \(-0{,}5\) e \(-0{,}8\). A
surpresa do primário quase não se correlaciona com as revisões (0,09 a 0,19).

\subsection{Prós e contras de cada medida}

\paragraph{Focus, resultado primário, 12 meses.}
Prós: cobre toda a amostra e preserva os 220 meses da base; é a medida mais próxima
do conceito da eq. (20) da nota, revisão da expectativa de superávit em evento fixo,
sem o efeito mecânico de rolar o horizonte; é reproduzível pela API Olinda; e é quase
a mesma informação do Prisma onde os dois coexistem. Contras: mede o setor público
consolidado, e não o Governo Central que emite a dívida estudada; 25\% dos meses têm
revisão nula, pela rigidez da mediana; usa \path{baseCalculo = 0}, mediana das
respostas dos últimos 30 dias, o que suaviza a revisão de fim de mês e espalha a
notícia entre meses vizinhos; e o horizonte de 12 meses é curto frente ao valor
presente dos superávits que a FTPL considera.

\paragraph{Focus, resultado primário, 36 meses.}
Prós: mais perto do respaldo de longo prazo; mesma cobertura; traz informação além do
12m (correlação de 0,52). Contras: mais ruído e mais zeros (30\%), porque poucos
respondentes revisam o terceiro ano; quase nenhuma relação com os mercados no mesmo
mês; sozinho, perde a notícia de curto prazo, mais ligada às necessidades de
financiamento.

\paragraph{Focus, DLSP, 12 e 36 meses.}
Prós: cobertura longa, quase sem zeros; resume primário, juros, crescimento e câmbio
num só número; é a variável usada por Carvalho Filho (2026). A correlação positiva com
a depreciação não é contaminação contábil, pois o setor público é credor líquido em
moeda estrangeira e a depreciação reduz a DLSP mecanicamente; é mais compatível com
notícia fiscal acompanhada de prêmio de risco. Contras: a expectativa de DLSP/PIB
também reage a Selic, inflação e PIB nominal, que são resultados ou canais das
projeções locais, de modo que purgar por ela pode condicionar em um mediador; não é
uma medida de respaldo; e o sinal frente ao respaldo é ambíguo, porque dívida
esperada maior pode vir de juros e não de déficit.

\paragraph{Prisma, resultado primário, 12 meses.}
Prós: conceito certo (Governo Central, o mesmo emissor); pesquisa mensal com
respondentes especializados; sem zeros; única medida cuja relação com os mercados tem
o sinal esperado (notícia boa com queda do EMBI e apreciação). Contras: começa em
2016 e corta a amostra para cerca de 110 meses, sem 2008 nem o período anterior à
crise de 2015; as edições não fecham no fim do mês e o calendário de coleta precisa
ser verificado; a extração ainda não é reproduzível.

\paragraph{Prisma, resultado primário, 36 meses.}
Prós: conceito certo e horizonte longo; única medida relacionada com a inflação
implícita (correlação de 0,23 com a variação da BEI de 5 anos). Contras: apenas 99
meses; o sinal da relação com a BEI é o oposto do esperado; mesmos problemas de data e
reprodutibilidade do 12m.

\paragraph{Surpresa do primário (realizado menos mediana Prisma), na data de divulgação.}
Prós: a mais próxima de notícia pura (AR(1) de 0,03); informação que de fato chega no
mês, independente de como os respondentes atualizam a mediana. Contras: mede o
resultado de um único mês, ruído contábil e de calendário, e não revisão do respaldo
futuro; data de divulgação aproximada (fim de \(t+1\)); começa em 2016. Funciona como
complemento das revisões, e não como substituto.

\paragraph{Séries do coautor descartadas.}
O Focus de resultado nominal inclui juros e tem endogeneidade maior que a DLSP; o
Focus de dívida bruta começa em 2018 e tem o mesmo problema de mediador; o Prisma de
dívida bruta muda de conceito em 2021 (Governo Geral até 2020, Governo Central
depois) e não pode ser emendado.

\subsection{A estratégia de Carvalho Filho (2026)}\label{sec:cf}

Carvalho Filho (2026), ``Fiscal Risk, Real Interest Rates and the Stock Market:
Evidence from Brazil'' (SSRN 7702078; pacote de replicação
doi:10.5281/zenodo.24206681, CC-BY-4.0), não usa uma revisão de pesquisa diretamente.
A notícia fiscal é um choque estrutural de risco fiscal num VAR bayesiano mensal
(2005-12 a 2026-08, 12 defasagens, prior Minnesota) identificado por restrições de
sinal:

\begin{itemize}[leftmargin=1.2cm]
    \item a variável fiscal é o nível da expectativa Focus de DLSP em horizonte fixo
    de 12 meses, média mensal das medianas, com o mesmo método de horizonte fixo que
    adaptamos em \path{build_fiscal_expectations_controls.py};
    \item as demais variáveis são IBC-Br, IPCA em 12 meses, Selic, juro real de 5
    anos, câmbio e Ibovespa real; VIX, Treasury de 10 anos e commodities entram como
    exógenas;
    \item o choque fiscal adverso eleva juntos, no impacto e no mês 1, a dívida
    esperada, o juro real de 5 anos e o dólar. O choque monetário se distingue por
    apreciar o câmbio e reduzir a inflação; o de demanda, por reduzir a dívida
    esperada;
    \item na robustez, a DLSP é trocada pela expectativa de déficit primário ou pela
    DLSP de 36 meses (correlação de 0,93 com o choque da base). A revisão de evento
    fixo aparece apenas na tabela narrativa dos oito episódios.
\end{itemize}

Das nossas opções, a mais próxima é a revisão Focus de DLSP de 12 meses. Com a série
de choques do pacote de replicação, as correlações são:

\begin{center}
\begin{tabular}{lc}
\toprule
Medida & Correlação com o choque de risco fiscal \\
\midrule
Focus DLSP 12m / 36m & 0,42 / 0,35 \\
Focus primário 12m / 36m & \(-0{,}09\) / \(-0{,}05\) \\
Prisma primário 12m / 36m & \(-0{,}09\) / \(-0{,}09\) \\
Surpresa do primário & 0,04 \\
\(u^N\) / \(u^F\) / \(u^R\) (base) & \(-0{,}09\) / 0,08 / \(-0{,}31\) \\
\bottomrule
\end{tabular}
\end{center}

As diferenças em relação à nossa DLSP são três: média mensal em vez de fim de mês;
resíduo de um VAR em vez de revisão bruta; e seleção, pelas restrições de sinal, da
parte da revisão que vem com alta do juro real e depreciação. Duas conclusões seguem.
Primeiro, revisões de primário e o choque de risco fiscal medem conceitos diferentes:
o respaldo esperado, como na eq. (20), e o risco fiscal precificado, caminho da dívida
mais prêmio. Segundo, a correlação de \(-0{,}31\) com \(u^R\) indica que, em meses de
choque fiscal adverso, o Tesouro emite menos NTN-B do que o previsto, provavelmente
recuando dos leilões indexados quando o juro real dispara. Esse componente é reação
da gestão da dívida às condições de mercado e faz sentido retirá-lo; mas, como o
choque é identificado com preços de ativos do próprio mês \(t\), usá-lo na purga traz
juro real e câmbio contemporâneos para dentro dos choques, com risco de viés nas
respostas de valuation. Além disso, o choque é uma estimativa com informação de toda a
amostra (regressor gerado). Por isso entra apenas como robustez.

\subsection{Purga: especificações}

O script \path{src/build_purged_shocks.py} estima, separadamente para cada classe,
\(u_t^j=a^j+\gamma^{j\prime}\Delta\mathcal{F}_t+\eta^{j\prime}W_t+\varepsilon_t^j\)
por MQO, com erros-padrão HAC, e grava
\path{data/clean/ftpl_monthly_panel_with_purged_shocks.csv}. Os choques recebem o nome \path{eps_{j}_{especificação de u}_{purga}}; o
choque principal é \path{eps_{j}_base_fprim}. As purgas, todas aplicadas a
\path{u_{j}_base}, são:

\begin{center}
\small
\begin{tabular}{p{0.17\textwidth}p{0.43\textwidth}p{0.3\textwidth}}
\toprule
Código & \(\Delta\mathcal{F}_t\) e tratamento & Amostra \\
\midrule
\path{fprim} & Focus primário 12m e 36m (\textbf{principal}) & 2006-11 a 2025-02 (220) \\
\path{fprim_xev} & idem, \(\gamma\) estimado sem os meses de evento (2007-04 e 2020-03 a 07) e aplicado a todos os meses & 220 (214 na estimação) \\
\path{fprim_dum} & idem, com as dummies de evento em \(W_t\); \(\varepsilon=0\) nesses meses & 220 \\
\path{fprim16} & Focus primário na amostra do Prisma, para separar medida de amostra & 2016-03 a 2025-02 (99) \\
\path{fprim_bc1} & Focus primário 12m e 36m com \path{baseCalculo = 1} & 2021-04 a 2025-02 (47) \\
\path{fprim21} & Focus primário principal na amostra de \path{fprim_bc1} & 47 \\
\path{pprim} & Prisma primário 12m e 36m & 2016-03 a 2025-02 (99) \\
\path{pprim_surp} & Prisma primário e surpresa do primário & 99 \\
\path{fdlsp} & Focus DLSP 12m e 36m & 220 \\
\path{fprim_fdlsp} & primário e DLSP juntos & 220 \\
\path{cfrisk} & choque de risco fiscal de Carvalho Filho (restrições de sinal) & 220 \\
\path{cfrisk_narr} & idem, com restrições narrativas & 220 \\
\path{fprim_cfrisk} & Focus primário e choque de Carvalho Filho & 220 \\
\bottomrule
\end{tabular}
\end{center}

A purga principal também é aplicada a todas as variantes da equação de previsão
(\path{long}, \path{w12}, \path{pca}, \path{principal}, \path{nolead},
\path{stock}, \path{liq}, \path{ramp} e \path{base_rt}). As séries de Carvalho Filho
são lidas do pacote de replicação guardado em
\path{data/raw/zenodo_24206681_fiscal_svar/} (fora do Git), com SHA-256 no arquivo de
auditoria, e acrescentadas ao painel como \path{cf_fiscal_risk} e
\path{cf_fiscal_risk_narr}.

\subsection{Resultados}

\begin{center}
\small
\begin{tabular}{lccc}
\toprule
Purga & \(R^2\) (\(N\) / \(F\) / \(R\)) & \(p\) do teste conjunto de \(\gamma\) & menor \(\rho(\varepsilon,\varepsilon^{\mathrm{fprim}})\) \\
\midrule
\path{fprim} & 0,003 / 0,001 / 0,006 & \(\ge 0{,}32\) & -- \\
\path{fprim_xev} & 0,002 / 0,008 / 0,006 & \(\ge 0{,}70\) & 0,98 \\
\path{fprim_dum} & 0,039 / 0,037 / 0,026 & \(\ge 0{,}70\) & 0,98 \\
\path{fprim16} & 0,001 / 0,002 / 0,013 & \(\ge 0{,}68\) & 0,999 \\
\path{fprim_bc1} & 0,017 / 0,017 / \textbf{0,155} & 0,003 em \(R\); \(\ge 0{,}39\) nas demais & 0,92 (\(R\)) \\
\path{fprim21} & 0,066 / 0,004 / 0,107 & \(\ge 0{,}24\) & 0,95 \\
\path{pprim} & 0,007 / 0,001 / 0,004 & 0,045 em \(N\); \(\ge 0{,}70\) nas demais & 0,99 \\
\path{pprim_surp} & 0,027 / 0,020 / 0,026 & 0,011 em \(N\) & 0,98 \\
\path{fdlsp} & 0,003 / 0,012 / 0,010 & \(\ge 0{,}15\) & 0,99 \\
\path{fprim_fdlsp} & 0,011 / 0,014 / 0,020 & \(\ge 0{,}09\) & 0,99 \\
\path{cfrisk} & 0,008 / 0,006 / \textbf{0,098} & \(<0{,}001\) em \(R\) & 0,95 (\(R\)) \\
\path{cfrisk_narr} & 0,022 / 0,021 / \textbf{0,084} & \(<0{,}001\) em \(R\); 0,049 em \(N\) & 0,95 (\(R\)) \\
\path{fprim_cfrisk} & 0,011 / 0,006 / \textbf{0,104} & \(<0{,}001\) em \(R\) & 0,95 (\(R\)) \\
\bottomrule
\end{tabular}
\end{center}

Na amostra base, as revisões de pesquisa, de primário ou de DLSP, Focus ou Prisma,
explicam no máximo 2,7\% da variância das inovações (\path{pprim_surp}); só na
subamostra curta de 2021 em diante o \(R^2\) chega a 6,6\% em \(N\)
(\path{fprim21}) e 15,5\% em \(R\) (\path{fprim_bc1}, Subseção~\ref{sec:bc1}). O teste conjunto de \(\gamma=0\) só rejeita a 5\%
nas purgas com Prisma na equação de \(N\) (\(p=0{,}045\) sem a surpresa e 0,011 com
ela), com coeficientes individuais insignificantes e amostra de 99 meses; a
evidência é frágil. Fora das purgas com o choque de Carvalho Filho e da versão
\path{fprim_bc1} em \(R\) (0,92 a 0,95), os choques purgados têm correlação de
pelo menos 0,98 com o principal. Com as dummies de evento, o \(R^2\) sobe para 3--4\%,
mas quem explica o aumento são as próprias dummies, que zeram os choques de 2020;
os coeficientes das revisões continuam insignificantes. A única medida com poder
relevante é o choque de Carvalho Filho na equação de \(R\): \(\gamma^R\) tem
estatística \(t\) de \(-5{,}8\) e o \(R^2\) chega a 10\%, sem efeito relevante em
\(N\) e \(F\). A correlação entre os choques principais é 0,24 entre \(N\) e \(R\) e
praticamente nula nos demais pares. A purga principal aplicada às variantes da
equação de previsão tem \(R^2\le 0{,}013\) em todas, e os choques correspondentes
reproduzem as correlações das inovações. Os resultados estão em
\path{data/diagnostics/purged_shocks_{gamma,fit,correlations,validation,audit}.csv},
\path{purged_shock_lists.json} e na figura \path{data/figures/purged_shocks.png}.

\subsection{Revisões Focus com baseCalculo = 1}\label{sec:bc1}

Com \path{baseCalculo = 0}, a mediana de fim de mês ainda inclui respostas de até 30
dias antes, de modo que parte da notícia de \(t\) só aparece na revisão de \(t+1\) e
parte da de \(t-1\) aparece em \(t\). A versão com \path{baseCalculo = 1} usa apenas
respostas dos últimos quatro dias úteis. A comparação nos 68 meses comuns
(2021-04 a 2026-11, \path{data/diagnostics/focus_basecalculo_comparison.csv} e
\path{data/figures/focus_basecalculo_comparison.png}) é:

\begin{center}
\footnotesize
\begin{tabular}{lcccc}
\toprule
 & Primário 12m & Primário 36m & DLSP 12m & DLSP 36m \\
\midrule
\(\rho\) dos níveis (fim de \(t-1\)) & 0,98 & 0,76 & 0,99 & 0,95 \\
\(\rho\) das revisões no mesmo mês & 0,72 & 0,41 & 0,52 & 0,32 \\
\(\rho\)(revisão bc0 em \(t\), bc1 em \(t-1\)) & 0,41 & 0,14 & 0,16 & 0,15 \\
desvio-padrão da revisão, bc0 / bc1 & 0,12 / 0,18 & 0,12 / 0,28 & 0,50 / 0,81 & 0,83 / 1,97 \\
AR(1) da revisão, bc0 / bc1 & 0,51 / 0,02 & 0,18 / \(-0{,}30\) & 0,00 / \(-0{,}34\) & \(-0{,}06\) / \(-0{,}40\) \\
meses com revisão nula, bc0 / bc1 & 13\% / 0\% & 6\% / 6\% & 0\% / 0\% & 6\% / 0\% \\
\bottomrule
\end{tabular}
\end{center}

No primário de 12 meses a nova versão faz o que se esperava: a revisão deixa de ser
persistente (AR(1) de 0,51 para 0,02), não tem meses nulos, e a revisão antiga em
\(t\) se correlaciona 0,41 com a nova em \(t-1\), sinal de que a janela de 30 dias
arrasta notícia de um mês para o seguinte. Nos demais indicadores a nova versão troca
suavização por ruído: com poucos respondentes, a mediana muda com a composição de
quem respondeu nos últimos dias, e as revisões ficam mais voláteis e com
autocorrelação negativa (\(-0{,}30\) a \(-0{,}40\)), típica de erro de medida que se
reverte.

Na purga (47 meses, 2021-04 a 2025-02), a versão nova explica 1,7\% de \(u^N\) e de
\(u^F\) e 15,5\% de \(u^R\), com a revisão de 36 meses significativa em \(R\)
(\(t=3{,}2\); teste conjunto \(p=0{,}003\)): superávit esperado maior vem com mais
emissão de NTN-B do que a prevista. A versão principal na mesma amostra
(\path{fprim21}) explica 10,7\% de \(u^R\), sem significância (\(p=0{,}24\)). Com
47 meses, dois regressores e uma revisão de 36 meses ruidosa, o resultado é frágil.
Os choques purgados com a versão nova têm correlação de 0,99 (\(N\), \(F\)) e 0,92
(\(R\)) com o choque principal. Conclusão: a versão principal continua sendo a de
\path{baseCalculo = 0}, a única que cobre a amostra; \path{fprim_bc1} entra como
robustez para o período recente, e o resultado em \(R\) deve ser reavaliado quando a
amostra da ETTJ for estendida após 2025-02.

\subsection{Leitura e decisões pendentes}\label{sec:purga_leitura}

\begin{enumerate}[leftmargin=1.2cm]
    \item A escolha entre as revisões de pesquisa quase não altera os choques. Pesam
    mais o tratamento de 2020 (dummies na purga ou na projeção local) e o corte da
    amostra em 2016 exigido pelo Prisma.
    \item A medida principal é o Focus primário de 12 e 36 meses: segue o conceito de
    respaldo da eq. (20), não usa preços de ativos e cobre toda a amostra. O Prisma
    (com \path{fprim16} como comparação na mesma amostra) é a robustez principal; a
    DLSP e o choque de Carvalho Filho são as alternativas de risco fiscal precificado.
    \item O choque de Carvalho Filho explica cerca de 10\% de \(u^R\), enquanto as
    pesquisas explicam perto de zero. Isso informa a decisão de identificação pendente
    com o coautor: purgar as inovações também da notícia fiscal precificada, e não só
    da esperada, e a proposta do coautor de usar a notícia fiscal como choque
    (Subseção~\ref{sec:decisoes}).
    \item Pendências técnicas: datar o calendário do Prisma e a divulgação do
    resultado realizado e obter as planilhas originais do Prisma. As revisões Focus
    com \path{baseCalculo = 1} já foram construídas (Subseção~\ref{sec:bc1}).
\end{enumerate}

\section{Projeções locais}\label{sec:lp}

Esta seção descreve a especificação comum a todas as projeções locais e os
resultados por MQO. \textbf{Os resultados por MQO estão viesados} pela reação da
emissão ao mercado do próprio mês (Seção~\ref{sec:ident}); as estimativas com
instrumento estão na Seção~\ref{sec:prop1}. Eles são mantidos aqui como ponto de
partida e porque o viés que revelam motivou a estratégia atual.

\subsection{Especificação}

O script \path{src/build_local_projections.py} estima a eq.~(\ref{eq:doc_local_projection})
para \(h=0,\ldots,24\), com os três choques purgados juntos:
\[
y_{t+h}-y_{t-1}=\alpha_h+\sum_{j\in\{N,F,R\}}\beta_h^j\,\varepsilon_t^j
+\Gamma_h'X_{t-1}+\mu_{m(t)}+v_{t+h}.
\]
\begin{itemize}[leftmargin=1.2cm]
    \item \textbf{Unidade do choque.} \(\varepsilon_t^j\) é multiplicado por 1000, de
    modo que uma unidade é uma emissão inesperada de 0,1\% do PIB de 12 meses
    defasado, o mesmo valor para as três classes. É a comparação a que se refere a
    ordenação \(\kappa^N>\kappa^F>\kappa^R\). Um desvio-padrão do choque vale 1,55
    unidade em \(N\), 1,50 em \(F\) e 0,66 em \(R\).
    \item \textbf{Resultados.} Nível de preços (\(100\log\) IPCA, de modo que
    \(\beta_h\) é a inflação acumulada), expectativa Focus de 12 meses, meta Selic,
    LTN de 1 ano, juros reais de 5 e 10 anos, inflação implícita de 2 e 5 anos
    (Subseção~\ref{sec:resultados_vars}),
    câmbio (\(100\log\) BRL/USD), EMBI e, para o canal de restrição de recursos,
    IBC-Br (\(100\log\)).
    \item \textbf{Controles \(X_{t-1}\)} (todos datados de \(t-1\) ou antes), com
    \(p=3\) defasagens: variação mensal do próprio resultado; os três choques; inflação
    mensal do IPCA; variação da meta Selic; crescimento do IBC-Br; depreciação
    cambial; nível da Selic, do EMBI e da expectativa Focus em \(t-1\); os três
    componentes principais da curva prefixada (médias de \(t-1\)); dummies de mês.
    Controles que repetem a defasagem do próprio resultado são excluídos, e
    dummies colineares são retiradas automaticamente.
    \item \textbf{Amostra.} \(t\) de 2007-02 a 2025-02 (217 meses em \(h=0\)); com
    resultados observados até 2026-10, a amostra só encolhe após \(h=20\), exceto
    na inflação implícita, que termina em 2025-01.
    \item \textbf{Inferência.} Newey--West com janela \(\max\{nw(n),h+1\}\). Para
    janelas de horizontes (0--6, 7--12, 13--24 e 0--24) a covariância conjunta dos
    coeficientes de choque entre horizontes é estimada com os escores empilhados na
    amostra comum da janela; como \(X_{t-1}\) não depende de \(h\), bastam as linhas
    de choque de \((X'X)^{-1}X'e_{t,h}\). Com ela se calculam: a diferença média
    \(\bar d=\frac{1}{|W|}\sum_{h\in W}(\beta_h^{j}-\beta_h^{k})\) e seu \(p\)
    unilateral para \(\bar d>0\); o teste sup-\(t\) de igualdade em todos os horizontes
    e a banda simultânea de 90\% (20\,000 simulações); a fração de horizontes com
    o sinal previsto. O Wald conjunto de igualdade também é gravado, mas
    rejeita em quase tudo, inclusive no IBC-Br, porque inverte uma covariância de
    até 25 horizontes quase colineares com cerca de 200 observações; o sup-\(t\) é
    o teste conjunto usado.
\end{itemize}

\begin{center}
\small
\begin{tabular}{p{0.24\textwidth}p{0.62\textwidth}}
\toprule
Especificação & Variação \\
\midrule
\path{base} & choques \path{eps_{j}_base_fprim}, \(p=3\) \\
\path{lags1}, \path{lags6} & \(p=1\) e \(p=6\) \\
\path{evdum} & dummies de 2007-04 e 2020-03 a 2020-07 em \(X\) \\
\path{pre2020} & \(t\le\) 2019-12 \\
\path{ordlast} & variações em \(t\) de câmbio, EMBI, juro real de 5 anos, LTN de 1
ano e expectativa Focus em \(X\): emissão ordenada por último no mês, e o impacto
dessas variáveis é zero por construção \\
\path{nopurge} & inovações \(u_t^j\) sem purga \\
\path{cfrisk} & purga \path{fprim_cfrisk} (Focus primário e choque de Carvalho Filho) \\
\path{pprim}, \path{fprim16} & purga Prisma e Focus na amostra do Prisma (2016+),
\(p=1\); o Prisma não tem janeiro e fevereiro \\
\path{rt} & equação de previsão em pseudo tempo real \\
\path{ramp} & choque real ampliado (NTN-B + NTN-C) no lugar de \(R\) \\
\path{iv_naive}, \path{iv_naive_w}, \path{iv_naive_sep}, \path{iv_pred} & MQ2E com as
surpresas de oferta como instrumentos (Seção~\ref{sec:prop1}) \\
\path{rf_naive} & forma reduzida: surpresas de oferta no lugar dos choques \\
\bottomrule
\end{tabular}
\end{center}

\subsection{Resultados por MQO}\label{sec:lp_res}

A tabela resume as respostas da especificação base (erros-padrão entre parênteses).

\begin{center}
\footnotesize
\begin{tabular}{llccc}
\toprule
Resultado & \(h\) & \(N\) & \(F\) & \(R\) \\
\midrule
Nível de preços (\%) & 6 & \(-0{,}01\) (0,05) & 0,13 (0,05) & 0,10 (0,14) \\
 & 12 & \(-0{,}05\) (0,07) & 0,13 (0,07) & 0,20 (0,23) \\
 & 24 & \(-0{,}17\) (0,10) & 0,09 (0,11) & 0,08 (0,26) \\
Expectativa IPCA 12m (p.p.) & 3 & \(-0{,}04\) (0,02) & 0,03 (0,03) & 0,03 (0,03) \\
 & 12 & 0,00 (0,03) & 0,04 (0,03) & 0,09 (0,08) \\
Meta Selic (p.p.) & 12 & \(-0{,}17\) (0,13) & 0,17 (0,10) & 0,39 (0,31) \\
 & 24 & \(-0{,}33\) (0,09) & 0,07 (0,16) & 0,51 (0,35) \\
Juro real 5 anos (p.p.) & 0 & \(-0{,}04\) (0,01) & 0,00 (0,02) & \(-0{,}12\) (0,03) \\
 & 18 & \(-0{,}14\) (0,05) & 0,00 (0,05) & 0,33 (0,12) \\
Câmbio (\%) & 0 & \(-0{,}35\) (0,14) & 0,31 (0,21) & \(-1{,}55\) (0,37) \\
EMBI (p.b.) & 0 & \(-1{,}7\) (1,2) & 0,4 (1,7) & \(-8{,}8\) (2,7) \\
\bottomrule
\end{tabular}
\end{center}

\begin{enumerate}[leftmargin=1.2cm]
    \item \textbf{A ordenação prevista não aparece.} No nível de preços, a diferença
    média \(\beta^N-\beta^R\) em 0--24 meses é \(-0{,}16\) (0,14) e
    \(\beta^N-\beta^F\) é \(-0{,}16\) (0,08); nenhum horizonte tem \(\beta^N>\beta^R\)
    nem \(N>F>R\). O choque de prefixados tende a reduzir preços, expectativas,
    Selic e juros; o de LFT eleva o nível de preços em cerca de 0,13\% entre 3 e 12
    meses (\(t\approx 2{,}5\)). O sup-\(t\) rejeita a igualdade \(N=F\) no nível de
    preços (\(p=0{,}002\)) e na expectativa Focus (\(p=0{,}005\)), mas com o sinal
    contrário ao previsto. \(N=R\) não é rejeitada no nível de preços (\(p=0{,}13\)).
    \item \textbf{O choque \(R\) é impreciso na inflação} (erro-padrão de 0,14 a 0,26
    no nível de preços), porque tem um terço da variância dos demais. Ele é
    acompanhado de alta da Selic e dos juros reais após 12 meses (0,33 p.p. no juro
    de 5 anos em \(h=18\)).
    \item \textbf{Impacto em \(h=0\) indica emissão oportunista.} No mês do choque
    \(R\), o juro real de 5 anos cai 0,12 p.p., o EMBI cai 9 p.b. e o real se
    aprecia 1,6\%; no mês do choque \(N\), a LTN de 1 ano cai 0,07 p.p. e a
    inflação implícita de 2 anos cai 0,04 p.p. O Tesouro parece emitir mais quando
    o mercado do mês está favorável, já que a equação de previsão só usa
    informação de \(t-1\). A Seção~\ref{sec:ident} detalha o problema e as
    propostas para tratá-lo.
    \item \textbf{Robustez.} O sinal negativo de \(\bar d^{NR}\) no nível de preços se
    mantém em todas as especificações; ele quase desaparece antes de 2020
    (\(-0{,}02\)) e é maior na amostra do Prisma (\(-0{,}47\), 87 meses), onde o
    choque \(R\) eleva o nível de preços em 0,6\% em 12--24 meses. Purga, número de
    defasagens, dummies e bloco real ampliado mudam pouco os resultados.
\end{enumerate}

Os resultados estão em
\path{data/diagnostics/local_projections_{irf,tests,fit,shock_scale,audit}.csv},
\path{local_projection_lists.json}, nas bandas sup-\(t\) de
\path{data/intermediate/local_projections/} e nas figuras
\path{data/figures/local_projections_base.png} (resultados principais),
\path{_base_extra.png}, \path{_differences.png} (diferenças com bandas pontual e
simultânea) e uma por especificação de robustez (\path{_evdum}, \path{_pre2020},
\path{_cfrisk}, \path{_nopurge}, \path{_ordlast}).

\section{Identificação dentro do mês}\label{sec:ident}

\paragraph{O problema.} A equação de previsão (Seção~\ref{sec:prev}) usa apenas
informação de \(t-1\). Isso garante que \(u_t^j\) não seja previsível com o passado,
mas não que seja exógeno às condições de mercado de \(t\). A emissão bruta de um
mês soma cerca de sete dias de leilão. Em cada leilão, o Tesouro escolhe o lote
ofertado e a taxa máxima que aceita, e o mercado escolhe quanto demanda a cada
taxa. Se as condições do mês mudam (uma melhora do apetite por risco, uma queda dos
juros), o Tesouro aumenta os lotes e aceita mais propostas, e o mercado compra
mais. Escrevendo o choque estimado como
\[
u_t^j=s_t^j+\lambda^j d_t,
\]
onde \(s_t^j\) é a decisão de oferta do Tesouro que não responde ao mês e \(d_t\)
é a notícia de mercado dentro do mês, a projeção local estima
\[
\mathrm{plim}\,\hat\beta_h^j=\beta_h^j+\lambda^j\,
\frac{\mathrm{Cov}(d_t,\,y_{t+h}-y_{t-1}\mid X_{t-1})}{\mathrm{Var}(u_t^j\mid X_{t-1})}.
\]
O viés tem o sinal da resposta de \(y\) à notícia de mercado multiplicado pelo
sinal de \(\lambda^j\).

\paragraph{Evidência.} Os fatos abaixo apontam para \(\lambda^j\neq 0\) em \(N\) e
\(R\), e não em \(F\):
\begin{enumerate}[leftmargin=1.2cm]
    \item \textbf{Correlação com preços do mesmo mês} (216 meses,
    \path{offer_surprises_exogeneity.csv}). \(\varepsilon_t^R\) tem correlação
    de \(-0{,}31\) com a variação do juro real de 5 anos, \(-0{,}27\) com a
    depreciação cambial e \(-0{,}22\) com a variação do EMBI; \(\varepsilon_t^N\),
    de \(-0{,}22\), \(-0{,}22\) e \(-0{,}17\), e de \(-0{,}21\) com a LTN de 1 ano;
    \(\varepsilon_t^F\) fica entre \(-0{,}01\) e 0,13.
    \item \textbf{Respostas de impacto.} Em \(h=0\), o choque \(R\) vem com juro real
    0,12 p.p. menor (\(t=-3{,}8\)), real 1,6\% mais apreciado (\(t=-4{,}2\)) e EMBI 9
    p.b. menor (\(t=-3{,}3\)); o choque \(N\), com LTN de 1 ano 0,07 p.p. menor
    (\(t=-4{,}1\)), inflação implícita de 2 anos 0,04 p.p. menor e expectativa Focus
    0,015 p.p. menor. Para variáveis de mercado com efeito imediato isso poderia
    ser o efeito do choque, mas uma emissão maior inesperada deveria \emph{elevar}
    os juros, não reduzi-los. O sinal é o de reação da emissão ao mercado.
    \item \textbf{Leilões restritivos e lotes conjuntos.} Chamamos de
    \emph{restritivo} o leilão em que o Tesouro vende o lote inteiro: nele, a oferta
    é a restrição ativa e a quantidade emitida é igual ao lote. No leilão \emph{não
    restritivo}, vende-se menos que o lote, que funciona apenas como teto não
    atingido. A unidade correta é o lote, não o título. O arquivo do Tesouro tem uma
    linha por título, e quando vários vencimentos dividem um lote conjunto o lote é
    repetido em cada linha. Um grupo de títulos do mesmo papel, no mesmo dia, com a
    mesma oferta é tratado como lote conjunto quando as vendas somam no máximo o
    lote. Se somam exatamente o lote, o caso é claro; se somam menos, é ambíguo e é
    tratado como conjunto, com uma variante que o trata como lotes separados. Se as
    vendas superam o lote, são lotes separados. Desde 2006-11 há 297 lotes conjuntos
    claros e 262 ambíguos em \(R\), 167 e 124 em \(F\), e 31 e 77 em \(N\). Contar por
    título subestima a aceitação: um lote conjunto de duas NTN-B vendido por inteiro
    aparecia como dois leilões não restritivos.

    A fração de leilões restritivos desde 2006-11 é (a contagem por título, usada
    numa versão anterior, está no Apêndice~\ref{app:hist}):
    \begin{center}
    \small
    \begin{tabular}{lccc}
    \toprule
     & \(N\) & \(F\) & \(R\) \\
    \midrule
    por título (incorreto) & 76\% & 35\% & 28\% \\
    por lote, ambíguos como conjuntos & 78\% & 65\% & 63\% \\
    por lote, ambíguos como separados & 77\% & 56\% & 50\% \\
    \bottomrule
    \end{tabular}
    \end{center}
    A mediana do vendido sobre o ofertado por lote é 1,00 nas três classes. Em \(R\),
    entre metade e dois terços dos leilões vendem o lote inteiro; a emissão de
    NTN-B não é predominantemente limitada pela demanda. Nos leilões não
    restritivos, a quantidade vendida também não é pura demanda: o Tesouro escolhe a
    taxa de corte, e vender menos que o lote pode ser falta de demanda ou recusa de
    propostas caras. Separar os dois casos exigiria o volume total de propostas por
    leilão, que não está no arquivo histórico (a verificar se o Tesouro o divulga
    nos resultados).
    \item \textbf{A própria oferta reage ao mês.} O valor mensal ofertado (lotes
    conjuntos contados uma vez) em \(R\) tem correlação de \(-0{,}14\) com a variação
    do juro real de 5 anos, de \(-0{,}30\) com a depreciação cambial e de \(-0{,}19\)
    com a variação do EMBI; em \(N\), de \(-0{,}21\), \(-0{,}26\) e \(-0{,}22\). Logo,
    trocar a emissão pelo valor ofertado no mês não resolve o problema: o tamanho
    dos lotes também responde ao mercado do mês.
    \item \textbf{Ser restritivo depende pouco do mês, exceto em \(N\).} Os leilões
    restritivos são aqueles em que uma variação do lote passa um para um para a
    emissão. A fração mensal de lotes restritivos tem as correlações abaixo com as
    variações de preços do mesmo mês:
    \begin{center}
    \small
    \begin{tabular}{lccc}
    \toprule
     & \(N\) & \(F\) & \(R\) \\
    \midrule
    \(\Delta\) juro real 5 anos & \(-0{,}11\) & 0,02 & 0,09 \\
    \(\Delta\) LTN 1 ano & \(-0{,}06\) & 0,09 & 0,09 \\
    depreciação cambial & \(-0{,}13\) & \(-0{,}01\) & 0,08 \\
    \(\Delta\) EMBI & \(-0{,}24\) & 0,08 & 0,07 \\
    \bottomrule
    \end{tabular}
    \end{center}
    Em \(N\), o lote é vendido inteiro com mais frequência em meses de mercado
    favorável (EMBI caindo). Em \(F\) e \(R\) as correlações são pequenas. Ainda
    assim, ser restritivo é um resultado do leilão e não deve ser usado para
    selecionar a amostra.
\end{enumerate}
Os números vêm de \path{src/diagnose_auction_takeup.py}
(\path{data/diagnostics/auction_takeup_{by_class,by_year,monthly_corr,binding,}}
\path{lot_vs_bond,lot_types}.csv).

\paragraph{Consequência para o teste.} Em \(N\), o viés vai contra a ordenação
prevista. O Tesouro vende mais prefixados quando juros nominais e inflação
implícita caem, de modo que \(\hat\beta^{N,\pi}\) tende a ser viesado para baixo. Em
\(R\), a direção é ambígua: a demanda por proteção contra inflação aumenta quando o
risco inflacionário sobe, o que viesaria \(\hat\beta^{R,\pi}\) para cima, mas os
meses de emissão alta de NTN-B são meses de real apreciado, o que empurra para
baixo. Por isso a falha da ordenação na Seção~\ref{sec:lp} não pode ser lida como
rejeição da teoria antes de tratar esse problema. A especificação \path{ordlast}
(variações financeiras de \(t\) em \(X\)) só reduz a diferença média \(N-R\) no nível
de preços de \(-0{,}16\) para \(-0{,}09\), mas ela supõe que a emissão não afeta
preços dentro do mês e por isso não resolve o problema.

\paragraph{Proposta 1: surpresas de oferta por leilão.}
A unidade passa a ser o leilão \(k\), no dia \(\tau_k\). A decisão de oferta do
Tesouro é o lote \(O_k\), anunciado no edital antes do leilão. (Na implementação,
\(k\) é o dia de leilão de cada papel, somando os lotes do dia; ver
Seção~\ref{sec:prop1}.) O choque é a parte do
lote que não se prevê com a informação disponível antes do edital:
\[
\tilde s_k=\frac{P_{k,\tau_k-1}\,\bigl(O_k-\mathbb{E}[O_k\mid\mathcal{I}_{\tau_k^e-1}]\bigr)}{PIB^{12m}_{t-1}},
\]
em que \(\tau_k^e\) é a data do edital, \(P_{k,\tau_k-1}\) é o preço do título na
véspera (o que evita a valoração pelo preço do próprio leilão) e a previsão usa o
lote do mesmo título no leilão anterior, a regra de lote padrão do trimestre e as
variações de preços entre o leilão anterior e a véspera do edital. O choque mensal
é \(\tilde s_t^j=\sum_{k\in t,\,j}\tilde s_k\). Ele entra na projeção local como
instrumento externo para \(\varepsilon_t^j\) (LP-IV) ou diretamente como choque.

\emph{Versão ponderada pela probabilidade de o leilão ser restritivo.} Uma surpresa
no lote só altera a emissão se o leilão for restritivo. Em leilões não restritivos o
lote não é atingido e a surpresa não tem efeito, o que enfraquece o instrumento,
sobretudo em \(F\) e \(R\). Restringir a amostra aos leilões restritivos seria um
erro: ser restritivo é um resultado do leilão, que em \(N\) se correlaciona com as
condições do mês (item 5 acima), e selecionar por ele reintroduziria a dependência
que se quer eliminar. A
alternativa é ponderar a surpresa pela probabilidade \emph{prevista} de o leilão ser
restritivo, estimada só com informação anterior ao edital:
\[
\tilde s_k^{w}=\hat\pi_k\,\tilde s_k,\qquad
\hat\pi_k=\Pr\bigl(\mbox{leilão }k\mbox{ restritivo}\mid\mathcal{I}_{\tau_k^e-1}\bigr),
\]
com \(\hat\pi_k\) dado, por exemplo, por um logit com o estado do leilão anterior do
mesmo título, a taxa média de aceitação recente do título, a classe, o prazo e as
variações de preços até a véspera do edital. O estado é persistente. A
probabilidade de um lote ser restritivo, dado o estado do lote anterior com a mesma
composição de títulos, é:
\begin{center}
\small
\begin{tabular}{lccc}
\toprule
 & \(N\) & \(F\) & \(R\) \\
\midrule
anterior restritivo & 0,83 & 0,79 & 0,75 \\
anterior não restritivo & 0,63 & 0,41 & 0,44 \\
\bottomrule
\end{tabular}
\end{center}
(Contando por título, a persistência parecia muito maior em \(F\) e \(R\), 0,10 e 0,08
após um leilão não restritivo, mas isso vinha dos lotes conjuntos.) Como \(\hat\pi_k\) é
função de informação predeterminada, a ponderação concentra o instrumento nos
leilões em que ele tem força sem condicionar no resultado do leilão. A
exogeneidade continua dependendo de \(\tilde s_k\); a ponderação só melhora a
relevância. A previsão de \(\hat\pi_k\) deve usar só o passado (estimação
recursiva), como a variante em pseudo tempo real da equação de previsão.
\begin{itemize}[leftmargin=1.2cm]
    \item \textbf{Fontes.} (i) Histórico de leilões do Tesouro
    (arquivo \path{historico-leiloes-todos.xls} em \path{data/raw/}), já no projeto, com oferta,
    quantidade vendida, taxa média, taxa de corte e financeiro por leilão e por
    volta, de 2000 a 2026; (ii) datas e conteúdo dos editais, nos comunicados de
    leilão publicados pelo Tesouro Nacional e no conjunto de dados de leilões do
    Tesouro Transparente (a confirmar se traz a data do edital); (iii) preços na
    véspera: preços e taxas diárias do Tesouro Direto (Tesouro Transparente, desde
    2002), a ETTJ ANBIMA diária, já integrada até 2025-01, e as taxas indicativas
    do mercado secundário da ANBIMA.
    \item \textbf{Viabilidade: alta.} A oferta por leilão já está no projeto e cobre
    toda a amostra; a primeira versão, com previsão pelo lote anterior e preços da
    ETTJ, pode ser feita com os dados atuais. O custo está em levantar as datas dos
    editais; enquanto isso, pode-se supor edital na véspera e testar a robustez.
    \item \textbf{Riscos.} Em \(R\), entre um terço e metade dos leilões não são
    restritivos, e nesses a variação do lote muda pouco a quantidade vendida; o
    instrumento pode ser fraco. A versão ponderada por \(\hat\pi_k\) mitiga, mas não
    elimina, esse risco. É
    preciso reportar o primeiro estágio das duas versões (estatística \(F\) robusta,
    Montiel Olea--Pflueger) e usar inferência robusta a instrumento fraco
    (Anderson--Rubin).
    O lote ainda pode reagir a informação dos dias anteriores ao edital, que só
    está parcialmente controlada. A soma mensal herda a frequência baixa das
    variáveis macro, então o ganho é na exogeneidade, não no número de
    observações.
\end{itemize}

\paragraph{Proposta 2: identificação de alta frequência no anúncio do lote.}
Em vez de prever o lote, mede-se a surpresa pela reação do mercado ao anúncio,
como na identificação de alta frequência da política monetária. Para cada edital,
\[
z_k=\Delta y^{\mathrm{janela}}_{k},
\]
a variação de juros na janela em torno da publicação do edital. Para separar choque
de oferta de notícia de demanda, usam-se restrições de sinal entre preços. Mais
oferta de \(R\) eleva o juro real e reduz a inflação implícita, enquanto notícia de
demanda por proteção eleva a inflação implícita. Mais oferta de \(N\) eleva o juro
nominal, com efeito pequeno no real. As surpresas são somadas no mês e usadas como
instrumento externo na projeção local.
\begin{itemize}[leftmargin=1.2cm]
    \item \textbf{Fontes.} Horário de publicação dos editais (comunicados do Tesouro);
    preços intradiários dos contratos futuros de DI e de cupom de IPCA (DAP) da B3,
    por meio de dados de mercado da B3 (pagos) ou de provedores como Bloomberg ou
    Refinitiv; com dados diários, taxas indicativas da ANBIMA e preços do Tesouro
    Direto, com janela de um dia.
    \item \textbf{Viabilidade: média a baixa.} Com janela intradiária, a identificação
    é a mais limpa, mas depende de dados pagos e de horários de edital que precisam
    ser coletados. Com janela diária, a viabilidade é alta, mas a janela de um dia
    mistura o edital com outras notícias, e os editais saem com frequência (cerca de
    sete dias de leilão por mês), o que dilui o sinal. A experiência com choques de
    oferta de títulos em outros mercados sugere reações pequenas a lotes de rotina;
    o poder deve vir de anúncios de mudança de estratégia (planos anuais, mudanças
    de lote, leilões cancelados ou extraordinários).
    \item \textbf{Riscos.} Instrumento fraco; o anúncio pode revelar informação do
    Tesouro sobre suas necessidades fiscais (efeito informacional, que é justamente
    o canal de respaldo que a purga quer separar), o que exige controlar notícias
    fiscais no mesmo dia.
\end{itemize}

\paragraph{Proposta 3: limites por ordenação.}
Sem dados novos, é possível delimitar o efeito. A especificação \path{ordlast}
atribui todo o comovimento do mês à reação da emissão (limite em que a emissão não
afeta preços dentro do mês). A especificação \path{base} atribui tudo ao efeito da
emissão (limite em que a emissão não reage ao mês). Uma família intermediária
fixa \(\lambda^j\) em valores plausíveis, por exemplo calibrados pela elasticidade
da taxa de aceitação à variação de preços nos leilões, e mostra como
\(\beta_h^{N}-\beta_h^{R}\) muda.
\begin{itemize}[leftmargin=1.2cm]
    \item \textbf{Fontes.} As já usadas.
    \item \textbf{Viabilidade: imediata,} mas é uma análise de sensibilidade, não uma
    identificação. Serve para dizer se a conclusão sobre a ordenação depende do
    tamanho do viés.
\end{itemize}

\paragraph{Alternativas descartadas como solução principal.} Um instrumento baseado
na demanda regulatória de fundos de pensão por NTN-B (normas do CMN e dados da
Previc ou da Abrapp), previsto na nota como robustez, desloca a demanda, não a
oferta. Ele move juros diretamente e viola a exclusão para as variáveis de
mercado. Uma base narrativa de decisões de emissão (Relatório Mensal da Dívida,
PAF, comunicados de leilões extraordinários) é útil para validar os choques, mas
tem poucos eventos, e muitos deles, como os leilões extraordinários, são reações a
estresse de mercado.

\paragraph{Recomendação.} (A Proposta~1 já foi implementada,
Seção~\ref{sec:prop1}; as Propostas~2 e 3 seguem pendentes.) Implementar primeiro a Proposta~1 com os dados disponíveis
(lote anterior como previsão e edital na véspera), nas versões simples e ponderada
por \(\hat\pi_k\), reportando o primeiro estágio por classe, e a Proposta~3 como
sensibilidade imediata. A Proposta~2 com dados diários
pode servir de validação; com dados intradiários, só se houver acesso a dados da B3
ou de um provedor.

\section{Surpresas de oferta e projeções locais com instrumento}\label{sec:prop1}

\paragraph{Construção.} O script \path{src/build_offer_surprises.py} implementa a
Proposta~1 só com dados já disponíveis.
\begin{itemize}[leftmargin=1.2cm]
    \item \textbf{Unidade.} As linhas por título são agregadas em lotes (conjuntos
    contados uma vez) e os lotes em dias de leilão de cada papel (LTN, NTN-F, LFT,
    NTN-B, NTN-B Principal): 2\,212 dias em \(N\), 882 em \(F\) e 686 em \(R\). O dia
    é a unidade porque a composição dos lotes conjuntos muda de um leilão para
    outro, o que impede comparar lote a lote.
    \item \textbf{Valor ofertado.} \(OV_k=\sum_{\ell\in k}O_\ell\,P^{ref}_\ell\), em que
    \(P^{ref}_\ell\) é a média, sobre os títulos do lote, do último preço unitário de
    leilão de cada título antes do dia \(k\); para títulos novos usa-se a média do
    papel. Nenhum preço do próprio dia entra.
    \item \textbf{Edital.} A data não está nos dados; supõe-se edital na véspera, de
    modo que a informação anterior ao edital vai até dois dias úteis antes do leilão.
    \item \textbf{Surpresas.} \path{naive}: \(OV_k-OV_{k^-}\), em relação ao dia de
    leilão anterior do mesmo papel. \path{pred}: \(OV_k-OV_{k^-}\exp(x_k'\hat b)\),
    com \(\log(OV_k/OV_{k^-})\) previsto por MQO recursivo por classe (reestimado
    a cada mês, janela crescente, a partir de 2006-02, quando começa a ETTJ diária)
    com as variações dos juros PRE 1 e 5 anos e IPCA 5 anos desde o leilão anterior,
    o estado e a aceitação do dia anterior, a aceitação média dos últimos quatro, o
    intervalo desde o leilão anterior e o prazo médio. As versões \path{_w}
    multiplicam a surpresa por \(\hat\pi_k\), de um logit recursivo com as mesmas
    variáveis. A variante \path{naive_sep} trata os grupos ambíguos como lotes
    separados.
    \item \textbf{Instrumentos mensais.} Soma no mês dividida pelo PIB de 12 meses
    defasado (mesma unidade de \(\varepsilon\)): \path{offsurp_{j}_{naive,pred,}}
    \path{naive_w,pred_w,naive_sep}, no painel
    \path{ftpl_monthly_panel_with_offer_surprises.csv}, entrada de
    \path{build_local_projections.py}.
\end{itemize}

\paragraph{Previsão do lote e de \(\hat\pi_k\).} A previsão recursiva quase não
melhora a ingênua: \(R^2\) fora da amostra de 0,03 em \(N\), \(-0{,}02\) em \(F\) e
0,10 em \(R\) em relação a \(OV_{k^-}\). As variações de preços antes do edital
explicam pouco da mudança do lote, o que é bom para a exogeneidade em relação ao
mercado dos dias anteriores, mas faz da versão \path{pred} apenas a \path{naive}
com ruído. O logit discrimina pouco no nível do dia (AUC de 0,62 em \(N\), 0,71 em
\(F\) e 0,61 em \(R\)), porque um dia só é restritivo se todos os seus lotes forem.

\paragraph{Primeiro estágio.} Regressão mensal de cada \(\varepsilon_t^j\) nos três
instrumentos com os controles da projeção local (2007-02 a 2025-02), com \(F\)
robusto de Sanderson--Windmeijer para cada endógena:
\begin{center}
\small
\begin{tabular}{lcccc}
\toprule
Instrumento & meses & \(N\) & \(F\) & \(R\) \\
\midrule
\path{naive} & 217 & 9,4 & 33,3 & 6,1 \\
\path{naive_sep} & 217 & 9,7 & 9,2 & 9,6 \\
\path{naive_w} & 185 & 9,5 & 46,4 & 10,4 \\
\path{pred} & 185 & 0,2 & 0,3 & 0,1 \\
\bottomrule
\end{tabular}
\end{center}
O primeiro estágio é diagonal: a surpresa de cada classe só move o choque da
própria classe (coeficientes de 0,84, 0,64 e 0,27 em \path{naive}, \(t\) de 4,1,
7,2 e 2,2). Em \(F\) o instrumento é forte; em \(N\) fica no limite usual de 10; em
\(R\) é fraco na versão principal e só chega a 10 na ponderada, com amostra mais
curta. A versão \path{pred} não serve como instrumento.

\paragraph{Exogeneidade.} A correlação com as variações de preços do mesmo mês cai em
relação a \(\varepsilon\), sem desaparecer:
\begin{center}
\small
\begin{tabular}{lcccc}
\toprule
 & \multicolumn{2}{c}{\(\Delta\) juro real 5 anos} & \multicolumn{2}{c}{depreciação cambial} \\
 & \(\varepsilon\) & \path{naive} & \(\varepsilon\) & \path{naive} \\
\midrule
\(N\) & \(-0{,}22\) & \(-0{,}10\) & \(-0{,}22\) & \(-0{,}19\) \\
\(F\) & \(-0{,}01\) & 0,10 & 0,13 & 0,21 \\
\(R\) & \(-0{,}31\) & \(-0{,}18\) & \(-0{,}27\) & \(-0{,}13\) \\
\bottomrule
\end{tabular}
\end{center}
Na versão ponderada, \(R\) cai para \(-0{,}08\) e \(-0{,}14\). O teste de antecipação
(\path{offer_surprises_anticipation.csv}) regride o valor vendido na surpresa por
faixa de \(\hat\pi_k\). Em \(R\), nos dias com \(\hat\pi_k<0{,}3\), a inclinação é
0,52 (\(t=5{,}4\)), mas esses dias ainda são restritivos em 32\% dos casos, porque o
logit discrimina pouco. Por isso o teste não separa antecipação da demanda de efeito
mecânico e é inconclusivo.

\paragraph{Projeções locais com instrumento.} As especificações \path{iv_naive}
(principal), \path{iv_naive_w}, \path{iv_naive_sep} e \path{iv_pred} estimam a
eq.~(\ref{eq:doc_local_projection}) por MQ2E, com os três choques instrumentados
pelas três surpresas e os mesmos controles; \path{rf_naive} é a forma reduzida, com
as surpresas no lugar dos choques. No nível de preços, por 0,1\% do PIB:
\begin{center}
\footnotesize
\begin{tabular}{lcccc}
\toprule
 & \(h\) & \(N\) & \(F\) & \(R\) \\
\midrule
\path{base} (MQO) & 12 & \(-0{,}05\) (0,07) & 0,13 (0,07) & 0,20 (0,23) \\
\path{iv_naive} & 6 & 0,07 (0,13) & 0,20 (0,06) & \(-0{,}11\) (0,36) \\
 & 12 & 0,09 (0,20) & 0,21 (0,08) & \(-0{,}27\) (0,49) \\
 & 24 & \(-0{,}03\) (0,27) & 0,11 (0,18) & 0,09 (0,59) \\
\path{iv_naive_sep} & 12 & 0,06 (0,20) & 0,21 (0,10) & \(-0{,}10\) (0,39) \\
\bottomrule
\end{tabular}
\end{center}
\begin{enumerate}[leftmargin=1.2cm]
    \item \textbf{O viés de simultaneidade some em \(N\) e \(F\).} Com o instrumento, o
    choque \(N\) deixa de vir com juros em queda no impacto (juro real de 5 anos:
    \(+0{,}04\), erro-padrão 0,06, contra \(-0{,}04\) por MQO), e as diferenças negativas
    e significativas de MQO desaparecem. A diferença média \(\beta^N-\beta^R\) no
    nível de preços em 0--24 meses passa de \(-0{,}16\) (0,14) para \(+0{,}16\) (0,40);
    \(\beta^N-\beta^F\), de \(-0{,}16\) (0,08) para \(-0{,}12\) (0,15). Na expectativa
    Focus, \(\beta^N-\beta^R\) passa de \(-0{,}07\) para \(+0{,}07\) (0,10).
    \item \textbf{O resultado mais robusto é o de \(F\).} Uma emissão inesperada de
    LFT de 0,1\% do PIB eleva o nível de preços em cerca de 0,2\% em 6 a 12 meses, em
    todas as versões com instrumento e na forma reduzida (0,13 em \(h=12\), \(t=2{,}3\)).
    \item \textbf{\(R\) continua sem identificação.} O instrumento é fraco (\(F=6\)) e, no
    impacto, o choque \(R\) instrumentado ainda vem com juro real em queda
    (\(-0{,}46\), erro-padrão 0,26) e real apreciado (\(-3{,}9\%\), erro-padrão 2,3); na
    forma reduzida, a surpresa de NTN-B tem impacto de \(-0{,}14\) p.p. no juro real
    (\(t=-3{,}1\)). O tamanho dos lotes de NTN-B parece acompanhar a demanda do mês. A
    versão ponderada, com primeiro estágio mais forte, dá respostas de \(R\)
    instáveis e grandes (2,3\% em \(h=12\)), com 182 meses.
    \item \textbf{A ordenação continua indeterminada.} Com erros-padrão de 0,3 a 0,4
    na diferença \(N-R\), nenhuma das especificações rejeita a igualdade. O teste
    deixou de apontar na direção contrária à teoria, mas não tem poder para
    confirmá-la.
\end{enumerate}
\begin{sloppypar}
Os resultados estão em \path{data/diagnostics/}, nos arquivos
\path{offer_surprises_{audit,forecast_fit,anticipation,first_stage,exogeneity}.csv};
nas surpresas por dia de leilão, em \path{offer_surprises_auctions.csv} (pasta
\path{data/intermediate/offer_surprises/}); na figura
\path{data/figures/offer_surprises.png}; e, para as projeções, nas especificações
\path{iv_*} e \path{rf_naive} de \path{local_projections_{irf,tests}.csv}, com
figuras \path{local_projections_{iv_naive,iv_naive_w,rf_naive}.png}.
\end{sloppypar}

\paragraph{Pendências.}\label{sec:prop1_pend} (i) Inferência robusta a instrumento fraco
(Anderson--Rubin por subvetor) para as diferenças entre classes; (ii) datas reais dos
editais; (iii) para \(R\), uma fonte de variação de oferta que não acompanhe a
demanda: a Proposta~2 em anúncios de mudança de estratégia ou uma versão no nível
do lote com \(\hat\pi\) mais discriminante.

%======================================================================
\part{Dados e pipeline (referência)}
%======================================================================

\section{Cadeia de scripts e reprodução}\label{sec:pipeline}

\subsection{Princípios}

\begin{itemize}[leftmargin=1.2cm]
    \item \textbf{Cada estágio lê o painel do estágio anterior e grava um painel
    novo} em \path{data/clean/}, sem alterar colunas existentes. O último painel da
    cadeia é sempre o mais completo; hoje é
    \path{ftpl_monthly_panel_with_offer_surprises.csv}.
    \item \textbf{Arquivos brutos nunca são sobrescritos.} Downloads novos são
    gravados com carimbo UTC (\path{*_AAAAMMDDTHHMMSSZ.csv}); entregas de terceiros
    são copiadas uma vez com o início do SHA-256 no nome. A opção \path{--offline}
    reutiliza o arquivo bruto mais recente.
    \item \textbf{Predeterminação.} Variáveis que entram como controle são datadas
    de \(t-1\) ou antes (sufixo \path{_lag1}); variáveis contemporâneas
    (\path{_rev_curr}, resultados de \(t\)) só entram na purga ou como resultado.
    \item \textbf{Validações interrompem a execução} quando falham (emendas,
    reconstruções, colunas repetidas), e cada estágio grava um arquivo de auditoria
    em \path{data/diagnostics/}.
\end{itemize}

\subsection{Estágios}

As dimensões incluem a coluna \path{month}; todos os painéis têm 564 meses
(1979-12 a 2026-11). As dimensões são as da execução atual e mudam quando um
estágio anterior ganha colunas; esta é a única tabela do documento que as registra.

\begin{center}
\footnotesize
\begin{longtable}{p{0.04\textwidth}p{0.36\textwidth}p{0.36\textwidth}p{0.12\textwidth}}
\toprule
 & Script (\path{src/}) & Painel gravado (\path{data/clean/ftpl_monthly_panel*}) & Colunas \\
\midrule
\endhead
0 & \path{notebooks/ftpl_build_data_notebook.ipynb} & \path{_with_pib_normalization} & 58 \\
1 & \path{build_structural_controls.py} & \path{_structural_controls} & 192 \\
2 & \path{import_anbima_ettj.py} & (bruto da ETTJ em \path{data/raw/anbima/}) & -- \\
3 & \path{build_anbima_yield_curve_controls.py} & \path{_with_yield_curve_controls} & 300 \\
4 & \path{build_fiscal_expectations_controls.py} \path{[--offline | --fetch P1,P2]} & \path{_with_fiscal_expectations} & 344 \\
5 & \path{build_prisma_fiscal_news.py [--offline]} & \path{_with_fiscal_news} & 384 \\
6 & \path{build_valuation_outcomes.py} & \path{_with_valuation_outcomes} & 402 \\
7 & \path{build_stock_history_sgs.py [--offline]} & \path{_with_stock_history} & 452 \\
8 & \path{build_maturity_paf_hist.py [--offline]} & \path{_with_maturity_paf_hist} & 545 \\
9 & \path{build_breakeven_outcomes.py} & \path{_with_breakeven_outcomes} & 564 \\
10 & \path{build_issuance_forecasts.py [--offline]} & \path{_with_issuance_innovations} & 683 \\
11 & \path{build_purged_shocks.py} & \path{_with_purged_shocks} & 749 \\
12 & \path{build_inflation_expectations.py [--offline]} & \path{_with_inflation_expectations} & 757 \\
13 & \path{build_offer_surprises.py} & \path{_with_offer_surprises} & 775 \\
14 & \path{build_local_projections.py} & (só diagnósticos e figuras) & -- \\
\bottomrule
\end{longtable}
\end{center}

Scripts fora da cadeia: \path{diagnose_auction_takeup.py} (diagnóstico de leilões
por lote, importado pelo estágio 13) e \path{fill_manual_templates_official_sources.py}
(preenchimento dos templates de PAF e reserva de liquidez, Apêndice~\ref{app:templates}).
Para alterar um estágio, basta reexecutar a partir dele; os estágios seguintes leem
o painel novo. O estágio 14 lê também as surpresas do 13, e o 13 usa as curvas
diárias do 3.

\subsection{Saídas}

\begin{itemize}[leftmargin=1.2cm]
    \item \path{data/raw/}: brutos (leilões, eDPF, SGS em \path{bcb/}, Focus em
    \path{focus/}, ANBIMA em \path{anbima/}, PAF e RMD em \path{paf_sources/} e
    \path{liquidity_sources/}, entrega Prisma do coautor);
    \item \path{data/intermediate/}: bases diárias e por leilão, bandas sup-\(t\);
    \item \path{data/clean/}: painéis mensais;
    \item \path{data/diagnostics/}: auditorias, validações, testes e listas de
    variáveis em JSON (\path{*_lists.json});
    \item \path{data/figures/}: figuras.
\end{itemize}

\subsection{Reprodução}

\begin{itemize}[leftmargin=1.2cm]
    \item Python 3 com \path{pandas}, \path{numpy}, \path{scipy}, \path{matplotlib},
    \path{requests}, \path{openpyxl} e \path{xlrd}. A econometria usa só
    \path{numpy}/\path{scipy}, sem \path{statsmodels}. Os gráficos usam o backend
    \path{Agg}.
    \item Fora do Git, a obter antes de rodar: o pacote de replicação de Carvalho
    Filho (Zenodo, doi:10.5281/zenodo.24206681) em
    \path{data/raw/zenodo_24206681_fiscal_svar/}; \path{DadosJoao/} (entregas
    originais do coautor, já copiadas com hash para \path{data/raw/});
    \path{backups/}.
    \item A API ANBIMA não é usada; credenciais, se um dia forem, vão em variáveis
    de ambiente (\path{.env} está no \path{.gitignore}).
    \item Execução completa sem internet: notebook com \path{REFRESH = False} e os
    scripts com \path{--offline}, na ordem da tabela. Com internet, os estágios 4,
    5, 7, 8, 10 e 12 baixam versões novas das séries.
    \item Quando chegar a extensão da ETTJ:
    \path{python3 src/import_anbima_ettj.py --source <arquivo>} e repetir a cadeia a
    partir do estágio 3.
\end{itemize}

\section{Fontes e construção dos dados}\label{sec:dados}

Cada subseção descreve um bloco do painel no estado atual: fonte, regra de
construção, cobertura e ressalvas. O número do estágio remete à tabela da
Seção~\ref{sec:pipeline}.

\subsection{Leilões, emissão, estoque e macro (estágio 0)}

O notebook \path{notebooks/ftpl_build_data_notebook.ipynb} (formato Colab, com
\path{REFRESH} para baixar ou reutilizar os brutos) baixa o histórico de leilões do
Tesouro (\path{data/raw/historico-leiloes-todos.xls}), a base de estoque da DPF
(\path{data/raw/estoquedpf.csv}, eDPF, só desde 2017-09) e séries do SGS/BCB e do
IPCA; mapeia instrumentos para \(N\), \(F\) e \(R\); agrega a emissão bruta mensal
por data de liquidação (vendas de mercado: VENDA e EXTRA VENDA, valor aceito
positivo); e normaliza emissões e vencimentos pelo PIB nominal de 12 meses
defasado. Gera também as figuras descritivas de \path{data/figures/} (emissão por
classe, composição, estoque, Selic e IPCA). A correção da leitura do SGS de
novembro de 2026 está no Apêndice~\ref{app:hist}.

\begin{center}
\small
\begin{longtable}{p{0.40\textwidth}p{0.13\textwidth}p{0.37\textwidth}}
\toprule
Arquivo em \path{data/clean/} & Linhas & Conteúdo \\
\midrule
\endhead
\path{auction_operations_clean.csv} & 19.604 & operações de leilão, 2000-01 a
2026-06 \\
\path{gross_issuance_wide.csv} & 318 & emissão bruta mensal por classe (também em
formato longo) \\
\path{macro_monthly.csv} & 564 & Selic, IBC-Br, câmbio, IPCA \\
\path{debt_stock_monthly.csv} & 104 & estoque por classe do eDPF, 2017-09 a 2026-04;
substituído pelas séries longas da Subseção~\ref{sec:estoque} \\
\path{scheduled_maturities_proxy.csv} & 59 & proxy de vencimentos do eDPF;
substituída pelo cronograma da Subseção~\ref{sec:venc} \\
\path{ftpl_monthly_panel_with_pib_normalization.csv} & 564 & painel inicial com
\path{g_{j}_pib}, \path{m_{j}_pib}, estoques e participações \\
\bottomrule
\end{longtable}
\end{center}

\subsection{Controles estruturais (estágio 1)}

\path{src/build_structural_controls.py} (``Etapa A'', instruções em
\path{docs/EtapaA.md}) cria os controles predeterminados de primeira geração:
defasagens e médias móveis de emissão, participações e estoques defasados,
calendário, dummies de eventos (crise de 2008, taper tantrum, recessão de
2015--16, greve dos caminhoneiros, Covid, eleições), dummies e valores de operações
extraordinárias (troca, recompra, venda e compra extraordinárias), gaps do PAF e
reserva de liquidez. PAF e reserva de liquidez vêm de templates manuais preenchidos
com fontes oficiais (Apêndice~\ref{app:templates}). A lista completa de variáveis e
diagnósticos está no Apêndice~\ref{app:etapaA}.

Vários desses controles foram depois substituídos por versões mais longas:
participações e estoques pelas séries desde 2000 (Subseção~\ref{sec:estoque}),
vencimentos e gaps do PAF pelo cronograma e pelos gaps desde 2003
(Subseção~\ref{sec:venc}). As versões originais continuam no painel, mas a equação
de previsão usa as novas.

\subsection{Curva de juros ANBIMA (estágios 2 e 3)}

A ``Etapa B'' cria controles predeterminados de curva de juros
que poderão entrar no vetor \(Z_{t-1}\) da equação de previsão de emissões. O
objetivo é construir medidas de nível, inclinação e curvatura para curvas
prefixadas e reais indexadas ao IPCA, além de inflação implícita e fatores obtidos
por PCA.

\subsubsection{Script e painel de entrada}

Foi criado o script \path{src/build_anbima_yield_curve_controls.py}, acompanhado do
notebook auxiliar \path{notebooks/build_anbima_yield_curve_controls.ipynb}. O script usa
\path{pathlib.Path}, cria os diretórios necessários, inspeciona o painel mensal
disponível e preserva o painel estrutural original.

O painel selecionado automaticamente foi:
\begin{center}
\path{data/clean/ftpl_monthly_panel_structural_controls.csv}.
\end{center}

O painel derivado é:
\begin{center}
\path{data/clean/ftpl_monthly_panel_with_yield_curve_controls.csv}.
\end{center}

\subsubsection{Fontes e situação da API}

O script foi preparado para usar a API ANBIMA Developers, sem credenciais
hardcoded, por meio de variáveis de ambiente como
\path{ANBIMA_ACCESS_TOKEN}, \path{ANBIMA_CLIENT_ID},
\path{ANBIMA_CLIENT_SECRET}, \path{ANBIMA_TOKEN_URL} e
\path{ANBIMA_API_BASE}. A documentação oficial consultada foi:
\begin{center}
\url{https://developers.anbima.com.br/pt/}\\
\url{https://developers.anbima.com.br/pt/documentacao/precos-indices/apis-de-precos/titulos-publicos/}
\end{center}

Na primeira execução não havia credenciais nem \path{ANBIMA_API_BASE}
configurado. A API não foi chamada, nenhum valor de curva foi inventado e as
colunas de curva ficaram vazias. A fonte efetivamente usada hoje é a ETTJ ANBIMA
enviada pelo coautor (subseção~\ref{sec:ettj}), lida pelo script como arquivo local
em \path{data/raw/anbima/}. A tentativa de API continua registrada em
\path{data/diagnostics/anbima_yield_curve_audit.csv}.

Como fallback adicional, foi criado o template:
\begin{center}
\path{data/raw/anbima/anbima_yield_curve_manual_template.csv}.
\end{center}

O template aceita as colunas \path{date}, \path{curve_type},
\path{vertex_du}, \path{yield_pct}, \path{source_document} e \path{notes}.
Os tipos padronizados são \path{PRE}, \path{IPCA} e \path{BEI}. A documentação
auxiliar está em \path{data/raw/anbima/README_anbima_manual_template.txt}.

\subsubsection{Bases intermediárias e transformação}

O script aceita fontes em CSV, XLS, XLSX, JSON e Parquet e tenta detectar
automaticamente colunas de data, tipo de curva, vértice e taxa. As curvas são
armazenadas em formato longo e convertidas para os vértices-alvo de 63, 126, 252,
504, 756, 1260, 1764 e 2520 dias úteis. A interpolação é linear por data e tipo de
curva e não extrapola além dos vértices observados.

As bases produzidas são:
\begin{itemize}[leftmargin=1.2cm]
    \item \path{data/intermediate/anbima/anbima_yield_curves_daily_long.csv};
    \item \path{data/intermediate/anbima/anbima_yield_curves_daily_wide.csv};
    \item \path{data/intermediate/anbima/anbima_yield_curve_features_daily.csv};
    \item \path{data/intermediate/anbima/anbima_yield_curve_features_monthly.csv};
    \item \path{data/intermediate/anbima/anbima_yield_curve_pca_factors_monthly.csv}.
\end{itemize}

Com a ETTJ, as bases diárias cobrem de 02/01/2006 a 31/01/2025 (4.786 dias úteis).
Os vértices-alvo coincidem com vértices da ETTJ, de modo que não foi necessária
nenhuma interpolação. Os meses sem curva continuam com \path{NaN}, sem valores
fictícios.

\subsubsection{Variáveis e predeterminação}

Para cada curva, o código está preparado para calcular os vértices
\path{3m}, \path{6m}, \path{1y}, \path{2y}, \path{3y}, \path{5y}, \path{7y} e
\path{10y}. A inflação implícita é calculada apenas quando as curvas prefixada e
IPCA existem na mesma data e no mesmo vértice:
\[
    BEI_h = PRE_h - IPCA_h.
\]

As medidas simples incluem nível de cinco anos, inclinações de 5--1 e 10--2 anos
e curvatura de 1--5--10 anos. Para respeitar a ordem temporal da equação de
emissões, foram criadas duas agregações mensais:

\begin{enumerate}[leftmargin=1.2cm]
    \item último dia útil observado até o fim do mês anterior, com sufixo
    \path{_eom_lag1};
    \item média das observações diárias do mês anterior, com sufixo
    \path{_avg_lag1}.
\end{enumerate}

O código não usa observações do próprio mês da emissão para esses controles.

\subsubsection{PCA e diagnósticos da Etapa B}

O script contém PCA separado para os blocos prefixado, IPCA e BEI, com até três
componentes, padronização das variáveis e opção de janela expansiva com janela
mínima de 60 meses. A implementação usa decomposição SVD do \path{numpy}, evitando
dependência obrigatória do \path{scikit-learn}.

Foram preparados os arquivos:
\begin{itemize}[leftmargin=1.2cm]
    \item \path{data/diagnostics/anbima_yield_curve_diagnostics.csv};
    \item \path{data/diagnostics/anbima_yield_curve_audit.csv};
    \item \path{data/diagnostics/anbima_yield_curve_descriptive_stats.csv};
    \item \path{data/diagnostics/anbima_yield_curve_correlations.csv};
    \item \path{data/diagnostics/anbima_yield_curve_pca_explained_variance.csv};
    \item \path{data/diagnostics/anbima_yield_curve_pca_loadings.csv};
    \item \path{data/diagnostics/yield_curve_control_lists.json}.
\end{itemize}

Na execução atual, as 72 variáveis de curva estão preenchidas de 2006-02 a 2025-02
(229 meses) e a lista de PCA contém 36 colunas (três componentes por bloco e por
agregação). Na amostra completa, o primeiro componente explica 96,7\% da variância
no bloco prefixado, 96,6\% no IPCA e 92,3\% na inflação implícita (versão
\path{_eom_lag1}). Os componentes e as medidas de nível, inclinação e curvatura usam
apenas vértices de 1 ano em diante.

\subsubsection{ETTJ ANBIMA enviada pelo coautor}\label{sec:ettj}

Em 22/11/2026 o coautor enviou \path{DadosJoao/data/data_yields.xlsx}, com a ETTJ
ANBIMA diária (curvas prefixada e real IPCA, inflação implícita e parâmetros de
Svensson), de 02/01/2006 a 31/01/2025, preparada para outro projeto. O arquivo
original da ANBIMA não será enviado, e o coautor vai estender a série para antes de
2006 e depois de janeiro de 2025.

A integração foi feita pelo novo script \path{src/import_anbima_ettj.py}:
\begin{enumerate}[leftmargin=1.2cm]
    \item a planilha é guardada intocada em
    \path{data/raw/anbima/original/data_yields__sha256_ebf1e33b96a6.xlsx}. O nome
    inclui o início do SHA-256; uma cópia idêntica é reaproveitada e uma versão
    diferente com o mesmo nome interrompe a execução, de modo que a extensão futura
    entra ao lado sem sobrescrever nada;
    \item o script lê apenas a cópia guardada. As 258 datas gravadas como texto
    (DD/MM/AAAA) são lidas com o dia primeiro;
    \item as curvas prefixada e IPCA são gravadas no formato do projeto (\path{date},
    \path{curve_type}, \path{vertex_du}, \path{yield_pct},
    \path{source_document}, \path{notes}), com taxas em \%, em
    \path{data/raw/anbima/anbima_ettj_daily_long.csv} (76.576 linhas). A inflação
    implícita não é gravada, porque a Etapa B a recalcula como \(PRE - IPCA\);
    \item os parâmetros diários de Svensson vão para
    \path{data/intermediate/anbima/anbima_ettj_svensson_params_daily.csv}.
\end{enumerate}

A validação, em \path{data/diagnostics/anbima_ettj_import_validation.csv},
interrompe o script em caso de falha. Todos os itens passaram: datas idênticas nas
quatro abas, estritamente crescentes e sem fins de semana; nenhum valor ausente; a
inflação implícita entregue é exatamente nominal menos real; as taxas coincidem com
a curva de Svensson reconstruída dos parâmetros, com \(t = du/252\) (erro máximo de
\(4\times10^{-15}\)); e o vértice de 1 ano difere, em média, 0,026 ponto percentual
da taxa ANBIMA da LTN de 1 ano (correlação de 0,9999 em 229 meses). O arquivo
\path{fatores_curva_juros.xlsx} da entrega anterior corresponde às linhas de fim de
mês desta série.

\subsubsection{Execução e resultado}

Após a importação, foram executados \path{src/build_anbima_yield_curve_controls.py}
e, em seguida, os estágios seguintes, em modo offline quando disponível:
expectativas Focus, notícia fiscal Prisma, variáveis de valuation e estoque desde
2000 (emenda com status PASS). A predeterminação foi conferida contra o dado bruto:
\path{pre_1y_eom_lag1} de fevereiro de 2006 é igual ao valor de 31/01/2006, e
\path{pre_1y_avg_lag1} é igual à média de janeiro.

Ressalvas:
\begin{itemize}[leftmargin=1.2cm]
    \item os vértices de 3 e 6 meses da curva real não são confiáveis: entre 2013 e
    2021, sem NTN-B curta, eles resultam da extrapolação da curva, com 172 valores
    fora do intervalo de \(-3\%\) a \(20\%\). As variáveis \path{ipca_3m},
    \path{ipca_6m}, \path{bei_3m} e \path{bei_6m} não devem ser usadas, e o
    importador marca esses vértices em \path{notes};
    \item os controles de curva só existem de 2006-02 a 2025-02, o que limita a
    amostra da equação de previsão (2006-11 a 2025-02) até que chegue a extensão do
    coautor (Seção~\ref{sec:pipeline}).
\end{itemize}
\subsection{Expectativas fiscais e notícia fiscal (estágios 4 e 5)}\label{sec:fiscal_dados}

\subsubsection{Expectativas Focus}

O script \path{src/build_fiscal_expectations_controls.py} constrói, a partir da API
Olinda do Banco Central, as expectativas Focus de dívida líquida do setor público e
de resultado primário, em \% do PIB, convertidas para horizontes fixos de 12 e 36
meses. As versões \path{_eom_lag1}, \path{_avg_lag1} e \path{_rev_lag1} são
predeterminadas; as revisões \path{_rev_curr} são notícia fiscal contemporânea e
destinam-se apenas à purga.

As séries principais usam \path{baseCalculo = 0}: a mediana das respostas mais
recentes informadas nos 30 dias anteriores à data de cálculo. O
script também baixa as medianas com \path{baseCalculo = 1}, que usam só as respostas
informadas a partir do quarto dia útil anterior, e grava as mesmas variáveis com o
prefixo \path{focus_primary_bc1} e \path{focus_dlsp_bc1} (16 colunas). A API só tem
essa versão a partir de 31/03/2021, com 8 a 20 respondentes por ano de referência,
contra 20 a 60 na versão principal; a primeira revisão é de 2021-04. Por isso ela é
robustez e não substitui a série principal (Subseção~\ref{sec:bc1}). O argumento
\path{--fetch} permite consultar a API só para alguns prefixos, mantendo os demais
nos arquivos brutos já salvos.

\subsubsection{Base Prisma do coautor}

A base de notícia fiscal construída pelo coautor João a partir do Prisma Fiscal e do
Focus foi integrada pelo script \path{src/build_prisma_fiscal_news.py}. A entrega é
copiada uma única vez, com hash SHA-256, para
\path{data/raw/prisma_fiscal_news_joao/}, e apenas essa cópia é lida. O painel
resultante é \path{data/clean/ftpl_monthly_panel_with_fiscal_news.csv}.

As principais variáveis são:

\begin{itemize}[leftmargin=1.2cm]
    \item \path{prisma_primary_12m_rev_curr_pib} e
    \path{prisma_primary_36m_rev_curr_pib}: revisão da mediana Prisma do resultado
    primário do Governo Central no mês \(t\), calculada por ano-alvo e combinada com
    os mesmos pesos de horizonte usados no Focus, em pontos percentuais do PIB
    acumulado em 12 meses defasado. São candidatas a \(\Delta\mathcal{F}_t^P\) e não
    devem entrar como controles da equação de emissão;
    \item \path{prisma_primary_12m_rev_lag1_pib} e
    \path{prisma_primary_36m_rev_lag1_pib}: versões predeterminadas;
    \item \path{prisma_primary_surp_pib_release}: resultado primário realizado menos
    a mediana Prisma para o mesmo mês, datada no mês aproximado de divulgação
    (\(t+1\)).
\end{itemize}

O sinal é positivo quando aumenta o superávit esperado, isto é, quando a notícia
fiscal é boa. As colunas originais do coautor são mantidas para rastreabilidade.

\subsubsection{Validação e problemas tratados}

As revisões do coautor para os horizontes \(h=0,\dots,3\) foram reproduzidas a partir
da base longa com diferença máxima da ordem de \(10^{-11}\). A revisão Prisma de 12
meses tem correlação de 0,95 com a revisão Focus de 12 meses. Dois problemas foram
tratados no código:

\begin{enumerate}[leftmargin=1.2cm]
    \item datas fora do primeiro dia do mês na série realizada (Apêndice~\ref{app:hist});
    \item o PIB usado na normalização é baixado diretamente da série SGS 4380, com o
    arquivo bruto preservado em \path{data/raw/bcb/}.
\end{enumerate}

Permanecem em aberto: a entrega das planilhas originais do Prisma e das consultas
Focus usadas pelo coautor; a quebra de conceito da
dívida bruta no Prisma (Governo Geral em 2015--2020 e Governo Central a partir de
2021), que impede a emenda das duas séries; a datação exata da divulgação do
resultado realizado; e a proposta do coautor de usar a notícia fiscal como choque
(Subseção~\ref{sec:decisoes}).

\subsection{Variáveis de resultado}\label{sec:resultados_vars}

\subsubsection{Juros, spread soberano e câmbio (estágio 6)}

O script \path{src/build_valuation_outcomes.py} adiciona variáveis de resultado a
partir dos arquivos mensais congelados do pacote de replicação de Carvalho Filho
(2026), doi:10.5281/zenodo.24206681, sob licença CC-BY-4.0. O pacote não é versionado
no Git e deve ser baixado do Zenodo para \path{data/raw/zenodo_24206681_fiscal_svar/}.


\begin{longtable}{p{0.30\textwidth}p{0.42\textwidth}p{0.18\textwidth}}
\toprule
\textbf{Variável} & \textbf{Origem} & \textbf{Cobertura} \\
\midrule
\endhead
\path{embi_spread_bp} & EMBI+ Brasil (Ipeadata) até julho de 2024; depois EMBI
Global Diversified Brasil mais o gap médio de 2,5 bp & 1994-04 a 2026-11 \\
\path{real_yield_5y_pct}, \path{real_yield_10y_pct} & NTN-B do Tesouro Direto,
interpolação diária a prazo constante & 2004-12 a 2026-11 \\
\path{ltn_1y_pct} & taxa ANBIMA da LTN de 1 ano (Ipeadata) & 2000-04 a 2026-11 \\
\path{cds5y_bp} & CDS soberano de 5 anos (Investing.com), robustez & 2009-01 a 2026-11 \\
\path{embi_gd_rel_bp} & EMBI GD Brasil menos EMBI GD Global, robustez & 2007-10 a 2026-11 \\
\path{swap_di_pre_360d_pct} & swap DI\(\times\)pré de 360 dias (B3), robustez &
1999-09 a 2025-11 \\
\bottomrule
\end{longtable}

Cada variável também entra na versão \path{_avg_lag1}, predeterminada. Todas as séries
são médias mensais de dados diários; por isso não há versão \path{_eom_lag1}. Não se
constrói inflação implícita neste estágio, pois o pacote não contém juros nominais e
reais no mesmo prazo; ela vem da ETTJ ANBIMA (Subseção~\ref{sec:bei}). Na emenda do EMBI, o EMBI+ e o EMBI GD têm correlação de 0,994 nos níveis no
período comum; o EMBI emendado e o CDS têm correlação de 0,92 nos níveis.

\subsubsection{Inflação implícita mensal (estágio 9)}\label{sec:bei}

O script \path{src/build_breakeven_outcomes.py} constrói variáveis de resultado de
inflação implícita a partir das curvas diárias da ETTJ. Para os prazos de 1, 2, 3, 5, 7 e 10
anos:
\begin{itemize}[leftmargin=1.2cm]
    \item \path{bei_{h}_pct}: média mensal em \(t\), a variável principal, comparável
    às demais variáveis de valuation;
    \item \path{bei_{h}_eom_pct}: último dia útil de \(t\);
    \item \path{bei_{h}_fisher_pct}: versão de Fisher, \((1+PRE)/(1+IPCA)-1\), para
    robustez.
\end{itemize}
A inflação implícita é a diferença simples \(PRE - IPCA\), como na Etapa B. Os prazos
de 3 e 6 meses foram excluídos. A validação confere que a diferença vale todos os
dias e que \path{bei_{h}_pct} em \(t\) é igual a \path{bei_{h}_avg_lag1} em \(t+1\) (e
o mesmo para o fim do mês), exatamente nos 229 meses. A cobertura vai de 2006-01 a
2025-01, até a chegada da extensão da ETTJ.

\subsubsection{Expectativa de inflação Focus (estágio 12)}\label{sec:focus_ipca}

\begin{sloppypar}
A expectativa de inflação é um dos resultados principais da nota. O script \path{src/build_inflation_expectations.py} baixa da API Olinda a
mediana diária do Focus para o IPCA dos próximos 12 meses
(série \path{ExpectativasMercadoInflacao12Meses}, com \path{baseCalculo = 0}) e grava um
arquivo bruto novo com carimbo UTC (\path{data/raw/focus/focus_ipca12m_*.csv}). As
variáveis são \path{focus_ipca12m_pct} (média mensal da série suavizada, a principal),
\path{focus_ipca12m_eom_pct} (último dia do mês),
\path{focus_ipca12m_nonsmooth_pct} (série não suavizada, correlação de 0,996 com a principal) e as versões
\path{_lag1}. A cobertura vai de 2001-12 a 2026-11; \path{--offline} lê o arquivo
bruto mais recente.
\end{sloppypar}

\subsection{Estoque e participações por classe desde 2000 (estágio 7)}\label{sec:estoque}

As colunas de estoque e participação por classe do painel vêm do arquivo
\path{estoquedpf.csv} do Tesouro (eDPF), que só começa em setembro de 2017. Como as
participações defasadas entram nas equações de emissão, isso restringia a amostra.
O script \path{src/build_stock_history_sgs.py} estende essas variáveis para
2000-02 a 2026-09 com séries do SGS/BCB de títulos federais em poder do mercado
(``posição em carteira'', excluída a carteira do Banco Central), em R\$ milhões.

\begin{longtable}{p{0.12\textwidth}p{0.55\textwidth}p{0.23\textwidth}}
\toprule
\textbf{SGS} & \textbf{Série} & \textbf{Uso} \\
\midrule
\endhead
10634 & LFT & \(F\) \\
10642 & NTN-B & \(R\) \\
10649 & LTN & parcela de \(N\) \\
10651 & total de prefixados & \(N\) \\
10650 & item prefixado residual, 2002-02 a 2005-09, até R\$ 0,12 bi & excluído \\
10647 & NTN-C & cobertura \\
4154 & total de títulos do Tesouro em mercado & denominador \\
\bottomrule
\end{longtable}

O bloco nominal é \(N = 10651 - 10650\), isto é, LTN mais NTN-F. A NTN-F não tem série
própria no SGS; o valor implícito nunca é negativo, a não ser por arredondamento de
\(\pm 1\), e passa a ser relevante em fevereiro de 2005. Os arquivos brutos são salvos
em \path{data/raw/bcb/} com carimbo de data e hora UTC e nunca são sobrescritos; a
opção \path{--offline} reutiliza o download mais recente.

\subsubsection{Emenda e validação}

A regra de emenda usa o eDPF onde ele existe (2017-09 a 2026-04) e o SGS nos demais
meses (2000-02 a 2017-08 e 2026-05 a 2026-09). A coluna \path{stock_hist_source}
registra a origem de cada mês. Nos 104 meses comuns, os níveis do SGS e do eDPF
diferem em no máximo 0,004\%, as participações em menos de 0,001 ponto percentual e
as variações mensais têm correlação 1,0 sem defasagem; a data do SGS no dia 1 corresponde,
portanto, ao estoque de fim daquele mês. O script interrompe a execução se a
diferença de nível passar de 0,5\% ou se a NTN-F implícita ficar negativa. Nos meses
do eDPF, as novas participações reproduzem exatamente as colunas \path{share_*}.

\subsubsection{Variáveis e cobertura}

Para a especificação base, são 32 colunas novas, sem alterar nenhuma coluna existente:
\path{stock_{N,F,R,total_NFR}_hist_brl}; \path{share_{g}_hist}, sobre \(N+F+R\), com a
mesma definição do painel; \path{share_{g}_hist_mkt}, sobre o total em mercado;
\path{coverage_NFR_over_market}; e as versões predeterminadas
\path{share_*_hist_lag1}, \path{share_*_hist_mkt_lag1}, \path{stock_*_hist_pib_lag1}
e \path{gross_{g}_over_lag_stock_hist}.

No início da amostra, \(N+F+R\) cobre só parte da dívida em mercado: 75\% em
2000-02, 70\% no mínimo de setembro de 2002, 82\% em janeiro de 2005 e 99\% hoje. O
restante é sobretudo NTN-C, títulos cambiais e LFT-A/B. A NTN-B é zero antes de
janeiro de 2002 e a NTN-F antes de fevereiro de 2005; essas participações são
estruturalmente nulas, não ausentes. A especificação base usa as participações
sobre \(N+F+R\); as participações sobre o total em mercado e a amostra a partir de
2005 ficam como robustez.

\subsubsection{Bloco real ampliado (robustez)}

A NTN-C é indexada ao IGP-M e, portanto, também é dívida real protegida da
inflação. Antes de 2006 ela dominava o bloco real: era 100\% do estoque real em
2001, 83\% em 2003 e 75\% em janeiro de 2005, contra 4\% hoje. Ela foi emitida em
leilões de 2001 a 2006, enquanto a NTN-B só passou a ser emitida com regularidade
em 2004 e em volume relevante em 2005. O mesmo script constrói, exclusivamente para
análise de robustez, o bloco \(R^{amp}\) = NTN-B + NTN-C. O bloco \(R\) da
especificação base continua sendo apenas a NTN-B.

A emissão bruta de NTN-C vem de \path{data/clean/auction_operations_clean.csv}, com
o mesmo filtro do notebook para a emissão base: vendas em mercado (VENDA e EXTRA
VENDA), valor aceito positivo e mês de liquidação. O script verifica que esse filtro
reproduz \path{gross_R_brl} (diferença máxima de R\$ 0,000004) antes de somar a
NTN-C. As vendas de NTN-C foram de R\$ 6,7 bi em 2001, 6,7 bi em 2002, 5,4 bi em 2003,
5,4 bi em 2004, 1,1 bi em 2005 e 1,5 bi em 2006. O estoque de NTN-C é a série 10647 do
SGS em todos os meses.

\begin{longtable}{p{0.09\textwidth}p{0.18\textwidth}p{0.18\textwidth}p{0.18\textwidth}p{0.17\textwidth}}
\toprule
\textbf{Ano} & \textbf{Meses com emissão \(R\)} & \textbf{Meses com emissão \(R^{amp}\)} &
\textbf{Emissão \(R^{amp}\) (R\$ bi)} & \textbf{Participação \(R^{amp}\)} \\
\midrule
\endhead
2000 & 0 & 0 & 0,0 & 1,4\% \\
2001 & 0 & 11 & 6,7 & 5,0\% \\
2002 & 0 & 9 & 6,7 & 9,9\% \\
2003 & 4 & 11 & 8,0 & 12,1\% \\
2004 & 11 & 12 & 9,7 & 14,0\% \\
2005 & 12 & 12 & 24,6 & 13,5\% \\
2006 & 12 & 12 & 51,2 & 21,4\% \\
\bottomrule
\end{longtable}

As 18 colunas adicionais são \path{gross_ntnc_brl}, \path{gross_Ramp_brl},
\path{gross_Ramp_brl_over_pib_lag1} (e \path{_pct}), \path{stock_ntnc_hist_brl},
\path{stock_Ramp_hist_brl}, \path{stock_total_NFRamp_hist_brl}; as participações
\path{share_N_hist_ramp}, \path{share_F_hist_ramp} e \path{share_Ramp_hist}, sobre
\(N+F+R^{amp}\), e \path{share_Ramp_hist_mkt}, sobre o total em mercado, com suas
versões \path{_lag1}; \path{stock_Ramp_hist_brl_lag1}, \path{stock_Ramp_hist_pib_lag1}
e \path{gross_Ramp_over_lag_stock_hist}. A lista
\path{robustness_broad_real_block_Ramp} em \path{stock_history_sgs_lists.json}
separa essas variáveis das da especificação base.

O uso previsto é um exercício de robustez com amostra longa, a partir de 2001, em
que o choque \(\varepsilon_t^{R^{amp}}\) substitui \(\varepsilon_t^R\) e as participações
ampliadas substituem as participações base. Duas ressalvas acompanham o exercício: a
NTN-C protege contra o IGP-M, não contra o IPCA; e, a partir de 2006, a diferença em
relação ao bloco base é pequena (de 1 a 6 pontos percentuais na participação, em
queda). Além disso, as metas do PAF para ``índices de preços'' incluem a NTN-C e se
referem à DPF total, enquanto \path{paf_gap_R_lag1} usa a participação apenas da
NTN-B sobre \(N+F+R\); essa diferença de definição fica registrada, mas não foi
alterada.

\subsection{Vencimentos programados e gaps do PAF desde 2003 (estágio 8)}\label{sec:venc}

A aproximação original de vencimentos (\path{scheduled_maturities_proxy.csv}) usa o
estoque do eDPF; por isso só existe a partir de 2017-10 e contém apenas o principal.
O script \path{src/build_maturity_paf_hist.py} constrói um cronograma mais longo e
recalcula os gaps do PAF com as participações longas. Ele lê
\path{data/clean/ftpl_monthly_panel_with_stock_history.csv} e grava
\path{data/clean/ftpl_monthly_panel_with_maturity_paf_hist.csv} (93 colunas novas).

\subsubsection{Fontes}

A fonte é o anexo do Relatório Mensal da Dívida (RMD) de maio de 2026, já preservado
em \path{data/raw/liquidity_sources/Anexo RMD_Maio_26.xlsx} e apenas lido (o SHA-256
fica na auditoria). O único download novo é o IGP-M (SGS 189), salvo em
\path{data/raw/bcb/sgs_189_<UTC>.csv}; a opção \path{--offline} reaproveita a cópia.

\begin{longtable}{p{0.10\textwidth}p{0.55\textwidth}p{0.25\textwidth}}
\toprule
\textbf{Aba} & \textbf{Conteúdo} & \textbf{Cobertura} \\
\midrule
\endhead
2.2 & estoque de fim de mês por título (quantidade, valor e vencimento) & 2010-01 a 2026-05 \\
1.3 & vencimentos realizados por título (principal e cupons) & 2006-11 a 2026-05 \\
3.3 & estoque a vencer em 12 meses, por indexador & 1999-12 a 2026-05 \\
2.4 & composição da DPF por indexador & 2000-12 a 2026-05 \\
3.7 & prazo médio da DPF & 2005-12 e 2006-11 a 2026-05 \\
\bottomrule
\end{longtable}

\subsubsection{Cronograma mensal}

A variável \path{maturity_{g}_sched_brl} mede o serviço da dívida devido em \(t\) e
conhecido no fim de \(t-1\), para \(g = N, F, R\) e \(C\) (NTN-C), além de
\path{Ramp} \(= R + C\) e \path{total_NFR}. As classes seguem o mapeamento do eDPF
(\(R\) inclui todas as variantes de NTN-B).
\begin{itemize}[leftmargin=1.2cm]
    \item principal: valor em \(t-1\) dos títulos que vencem em \(t\), mesmo conceito
    da aproximação do eDPF;
    \item cupons: quantidade em \(t-1\) vezes o cupom por título. NTN-F paga 10\% ao ano
    em janeiro e julho; NTN-B paga 6\% e NTN-C 6\% (12\% no título de 2031) sobre o
    VNA, no mês do vencimento e seis meses depois. O VNA usa o IPCA ou o IGP-M de
    \(t-2\), para que o valor seja predeterminado; o erro é de um mês de inflação
    sobre o cupom. NTN-B Principal não paga cupom, e as parcelas de NTN-B1
    (Renda+ e Educa+) não são modeladas.
\end{itemize}

O cronograma da aba 2.2 cobre \(t\) de 2010-02 a 2026-06. De 2006-11 a 2010-01 a
coluna usa os vencimentos realizados da aba 1.3; \path{maturity_sched_source} indica a
origem. Antes de 2006-11 não há cronograma mensal. Principal e cupons também ficam
separados, e os realizados ficam em \path{maturity_{g}_realized_rmd_brl}, que é
contemporâneo e serve apenas para validação.

A validação interrompe o script em caso de falha. O principal é idêntico ao da
aproximação do eDPF nos 59 meses em comum. Contra os vencimentos realizados, de
2010-02 a 2026-05, a diferença absoluta média em relação à média realizada é de
0,9\% em \(N\) e \(F\), 2,9\% em \(R\), 1,6\% em \(C\) e 1,2\% em \(N+F+R\), com
correlações de pelo menos 0,9997.

\subsubsection{Janelas de 12 meses, composição e prazo médio}

A partir da aba 3.3, com o estoque no fim de \(t-1\), são criadas
\path{maturity12m_{N,F,Ramp,total}_rmd_brl_lag1} e as participações correspondentes
(\path{_share_lag1}), desde 2000-01, além de \path{maturity12m_total_rmd_pib_lag1} e
\path{liquidity_reserve_over_maturity12m_rmd_lag1}. O bloco de índices de preços
inclui NTN-B e NTN-C. Das abas 2.4 e 3.7 vêm
\path{share_{pre,price,float,fx}_dpf_rmd_lag1}, \path{stock_dpf_rmd_brl_lag1} e
\path{avg_maturity_dpf_years_rmd_lag1}.

\subsubsection{Gaps do PAF}

Os gaps (\path{paf_gap_}, distâncias aos limites e indicadores de fora da banda) são
calculados contra as bandas já presentes no painel para três definições de
participação:
\begin{itemize}[leftmargin=1.2cm]
    \item \path{hist}: participações sobre \(N+F+R\) (\path{share_{g}_hist_lag1});
    \item \path{hist_mkt}: participações sobre o total em mercado;
    \item \path{rmd}: composição da DPF da aba 2.4, que é a métrica em que as bandas do
    PAF são escritas (inclui dívida externa e NTN-C), também para câmbio.
\end{itemize}
Os gaps passam a cobrir 2003-01 em diante; os antigos começavam em 2017-10, e
\path{paf_gap_R_hist_lag1} reproduz exatamente \path{paf_gap_R_lag1} onde ambos
existem. Foram criados também \path{paf_gap_maturing_12m_rmd_lag1}, com a parcela da
DPMFi a vencer em 12 meses, e \path{paf_gap_avg_maturity_rmd_lag1}, desde 2006-01.

A versão recomendada é a \path{rmd}. Com \(N+F+R\) no denominador, os gaps ficam muito
viesados antes de 2010, pois a dívida cambial e a NTN-C ficam de fora: a participação
Selic fica acima da banda em 71\% dos meses, e o gap médio da Selic em 2003 é de 50
pontos percentuais. As versões \path{hist} e \path{hist_mkt} ficam como robustez.

Saídas adicionais: \path{data/intermediate/maturity_paf/}, os diagnósticos
\path{maturity_paf_hist_{audit,validation,diagnostics}.csv},
\path{maturity_paf_hist_lists.json} e a figura
\path{data/figures/maturity_paf_hist.png}.

\section{Convenções e glossário}\label{sec:convencoes}

\subsection{Classes de dívida}

\begin{center}
\small
\begin{tabular}{p{0.14\textwidth}p{0.42\textwidth}p{0.32\textwidth}}
\toprule
Classe & Títulos & Observação \\
\midrule
\(N\) & LTN e NTN-F (prefixados) & \(\kappa^N=1\) \\
\(F\) & LFT (Selic) & \(0<\kappa^F<1\) \\
\(R\) & NTN-B, inclusive NTN-B Principal (IPCA) & \(\kappa^R=0\) \\
\(R^{amp}\) (\path{Ramp}) & NTN-B e NTN-C (IGP-M) & só robustez \\
\(C\) & NTN-C & usada no cronograma de vencimentos \\
\bottomrule
\end{tabular}
\end{center}

\subsection{Sufixos e unidades}

\begin{center}
\small
\begin{longtable}{p{0.26\textwidth}p{0.64\textwidth}}
\toprule
Sufixo ou prefixo & Significado \\
\midrule
\endhead
\path{_brl}, \path{_pib} & reais correntes; fração do PIB acumulado em 12 meses
defasado (\path{pib_12m_lag1_brl}) \\
\path{_pct}, \path{_bp} & percentual; pontos-base \\
\path{_lag1}, \path{_lag2} & defasagem de 1 ou 2 meses; predeterminado \\
\path{_eom}, \path{_avg} & último dia útil do mês; média do mês \\
\path{_eom_lag1}, \path{_avg_lag1} & último dia útil ou média do mês anterior;
controles de curva e expectativas \\
\path{_rev_curr} & revisão de expectativa no mês \(t\): notícia contemporânea, só
para a purga \\
\path{_rev_lag1} & revisão do mês anterior, predeterminada \\
\path{_hist} & séries longas de estoque (SGS emendado ao eDPF) \\
\path{_rmd} & calculado com o anexo do Relatório Mensal da Dívida \\
\path{_sched} & serviço da dívida programado, conhecido em \(t-1\) \\
\path{_bc1} & Focus com \path{baseCalculo = 1} (padrão: 0) \\
\path{u_{j}_{var}} & inovação da equação de previsão, variante \path{var} \\
\path{eps_{j}_{var}_{purga}} & choque purgado; principal: \path{eps_{j}_base_fprim} \\
\path{offsurp_{j}_{versão}} & surpresa de oferta mensal; principal:
\path{offsurp_{j}_naive} \\
\bottomrule
\end{longtable}
\end{center}

\paragraph{Sinais.} Revisões de primário: positivo é superávit esperado maior
(notícia boa). Revisões de DLSP e choque de Carvalho Filho: positivo é adverso.
Câmbio: alta é real mais fraco. Choques: uma unidade é emissão inesperada de 0,1\%
do PIB de 12 meses defasado.

\subsection{Notação}

\begin{center}
\small
\begin{longtable}{p{0.18\textwidth}p{0.72\textwidth}}
\toprule
Símbolo & Significado \\
\midrule
\endhead
\(\kappa^j\) & elasticidade do valor real da classe \(j\) ao nível de preços \\
\(\rho^j\) & parcela do choque respaldada por superávits futuros \\
\(g_t^j\) & emissão bruta da classe \(j\) sobre o PIB de 12 meses defasado \\
\(m_{t|t-1}^j\) & serviço da dívida devido em \(t\), conhecido no fim de \(t-1\) \\
\(Z_{t-1}\) & informação predeterminada da equação de previsão \\
\(u_t^j\) & inovação de emissão (resíduo da eq.~\ref{eq:doc_forecast}) \\
\(\Delta\mathcal{F}_t\) & notícia fiscal do mês (revisão da expectativa de superávit) \\
\(\varepsilon_t^j\) & choque purgado (resíduo da eq.~\ref{eq:doc_survey_purge}) \\
\(\beta_h^j\) & resposta em \(h\) meses a uma unidade de \(\varepsilon^j\) \\
\(X_{t-1}\) & controles da projeção local \\
\(\bar d^{jk}\) & diferença média \(\beta_h^j-\beta_h^k\) numa janela de horizontes \\
\(O_k\), \(OV_k\) & lote ofertado e valor ofertado no dia de leilão \(k\) \\
\(\hat\pi_k\) & probabilidade prevista de o leilão ser restritivo \\
leilão restritivo & leilão que vende o lote inteiro (a oferta é a restrição ativa) \\
lote conjunto & lote único dividido entre vários vencimentos do mesmo papel, repetido
em cada linha do arquivo do Tesouro \\
\bottomrule
\end{longtable}
\end{center}

%======================================================================
\appendix
%======================================================================

\section{Detalhes da Etapa A}\label{app:etapaA}

Registro da implementação original da Etapa A (junho de 2026), mantido para
consulta das listas de variáveis. Três ressalvas valem para todo o apêndice: (i) o
painel tinha então 559 meses (até 2026-06), e hoje tem 564; (ii) os templates de
PAF e de reserva de liquidez, descritos aqui como vazios, foram preenchidos em
seguida (Apêndice~\ref{app:templates}); (iii) as participações, estoques,
vencimentos e gaps do PAF desta etapa foram substituídos pelas versões longas das
Subseções~\ref{sec:estoque} e~\ref{sec:venc}, e a baixa cobertura das
participações apontada nos alertas deixou de existir.

\subsection{Implementação}

\subsubsection{Motivação}

A Etapa A implementa variáveis estruturais que devem entrar fora de qualquer PCA na
equação de previsão das emissões. A lógica é construir controles predeterminados,
isto é, informações disponíveis antes da decisão de emissão no mês \(t\). Esses
controles serão usados posteriormente para obter inovações de emissão \(u_t^j\).

\subsubsection{Arquivo de instruções}

As instruções estão em \path{EtapaA.md}. Elas especificam que o código deve:

\begin{itemize}[leftmargin=1.2cm]
    \item inspecionar a estrutura do diretório;
    \item identificar notebook, scripts e CSVs existentes;
    \item localizar o painel mensal mais recente;
    \item ler os nomes reais das colunas antes de programar;
    \item não sobrescrever arquivos originais;
    \item criar uma etapa nova do pipeline com script, notebook auxiliar, CSV final,
    templates manuais, diagnósticos e gráficos.
\end{itemize}

\subsubsection{Script implementado}

Foi criado o script \path{src/build_structural_controls.py}. O script:

\begin{itemize}[leftmargin=1.2cm]
    \item usa \path{pathlib.Path};
    \item detecta automaticamente o painel mais recente;
    \item imprime os nomes reais das colunas dos arquivos usados;
    \item cria \path{data/raw/}, \path{data/clean/},
    \path{data/diagnostics/} e \path{data/figures/};
    \item cria \path{panel_structural = panel.copy()} antes de transformar os dados;
    \item salva um arquivo derivado sem sobrescrever o painel original;
    \item é modular, com funções para diretórios, templates, lags, operações
    extraordinárias, PAF, reserva de liquidez, listas de controles e diagnósticos;
    \item força backend gráfico não interativo \path{Agg}, para funcionar localmente
    no macOS sem janela gráfica;
    \item usa apenas \path{pandas}, \path{numpy}, \path{matplotlib} e bibliotecas
    padrão na execução atual.
\end{itemize}

\subsubsection{Notebook auxiliar}

Foi criado \path{notebooks/build_structural_controls.ipynb}, um wrapper simples que executa
\path{src/build_structural_controls.py}. O objetivo é permitir uso em notebook ou Colab
sem duplicar a lógica do script.

\subsubsection{Painel usado como entrada}

O script selecionou como entrada:
\begin{center}
\path{data/clean/ftpl_monthly_panel_with_pib_normalization.csv}.
\end{center}
Esse painel contém 559 linhas e 58 colunas, cobrindo o período de dezembro de 1979
a junho de 2026. As colunas reais inspecionadas incluem:

\begin{itemize}[leftmargin=1.2cm]
    \item emissões: \path{g_N_pib}, \path{g_F_pib}, \path{g_R_pib},
    \path{g_total_NFR_pib};
    \item participações: \path{share_N}, \path{share_F}, \path{share_R};
    \item estoques: \path{stock_N_brl}, \path{stock_F_brl}, \path{stock_R_brl},
    \path{stock_total_NFR_brl};
    \item vencimentos: \path{m_N_pib}, \path{m_F_pib}, \path{m_R_pib};
    \item PIB: \path{pib_12m_lag1_brl};
    \item macro: \path{selic_target_pa}, \path{ibc_br_sa}, \path{brl_usd},
    \path{ipca_index}, \path{ipca_mom_pct}, \path{ipca_yoy_pct}.
\end{itemize}

\subsection{Variáveis construídas na Etapa A}

\subsubsection{Emissões defasadas}

Foram criadas defasagens configuráveis de emissões normalizadas por PIB, com
\path{N_LAGS_ISSUANCE = 3}. Para cada uma das séries disponíveis
\path{g_N_pib}, \path{g_F_pib}, \path{g_R_pib} e
\path{g_total_NFR_pib}, foram criadas defasagens \path{lag1}, \path{lag2} e
\path{lag3}.

Também foram criadas médias móveis predeterminadas:

\begin{itemize}[leftmargin=1.2cm]
    \item \path{g_N_pib_ma3_lag1};
    \item \path{g_F_pib_ma3_lag1};
    \item \path{g_R_pib_ma3_lag1}.
\end{itemize}

Essas médias usam apenas informação até \(t-1\).

\subsubsection{Participações e estoques defasados}

A partir de \path{share_N}, \path{share_F} e \path{share_R}, foram criadas:

\begin{itemize}[leftmargin=1.2cm]
    \item \path{share_N_lag1};
    \item \path{share_F_lag1};
    \item \path{share_R_lag1}.
\end{itemize}

O script também criou o diagnóstico \path{share_sum_structural_check} e o erro
\path{share_sum_abs_error}. Na execução local, o erro absoluto máximo da soma das
participações em relação a 1 foi \(2.22 \times 10^{-16}\), isto é, apenas erro de
arredondamento numérico.

Para estoques, foram criadas as defasagens de \path{stock_N_brl},
\path{stock_F_brl}, \path{stock_R_brl} e \path{stock_total_NFR_brl}. Como
\path{pib_12m_lag1_brl} existe, também foram criadas normalizações por PIB:

\begin{itemize}[leftmargin=1.2cm]
    \item \path{stock_N_pib_lag1};
    \item \path{stock_F_pib_lag1};
    \item \path{stock_R_pib_lag1};
    \item \path{stock_total_NFR_pib_lag1}.
\end{itemize}

\subsubsection{Vencimentos e janelas futuras}

A partir das variáveis \path{m_N_pib}, \path{m_F_pib} e \path{m_R_pib}, o script
criou \path{m_total_NFR_pib} quando necessário e depois construiu:

\begin{itemize}[leftmargin=1.2cm]
    \item \path{m_total_NFR_pib_lag0};
    \item \path{m_total_NFR_pib_lead1};
    \item \path{m_total_NFR_pib_window_3m};
    \item \path{m_total_NFR_pib_window_6m};
    \item \path{m_total_NFR_pib_window_12m};
    \item \path{m_N_pib_window_3m};
    \item \path{m_F_pib_window_3m};
    \item \path{m_R_pib_window_3m}.
\end{itemize}

O código contém comentário explícito de que janelas futuras de vencimentos só são
interpretáveis como predeterminadas se o cronograma futuro de vencimentos era
conhecido em \(t-1\), como ocorre com cronogramas contratuais.

\subsubsection{Calendário e eventos}

Foram criadas:

\begin{itemize}[leftmargin=1.2cm]
    \item \path{year};
    \item \path{quarter};
    \item \path{calendar_month}.
\end{itemize}

Também foram criadas dummies configuráveis:

\begin{longtable}{p{0.38\textwidth}p{0.42\textwidth}}
\toprule
\textbf{Dummy} & \textbf{Período} \\
\midrule
\endhead
\path{dummy_global_financial_crisis} & setembro de 2008 a junho de 2009 \\
\path{dummy_taper_tantrum} & maio de 2013 a dezembro de 2013 \\
\path{dummy_domestic_recession_2015_2016} & janeiro de 2015 a dezembro de 2016 \\
\path{dummy_truckers_strike} & maio de 2018 \\
\path{dummy_covid} & março de 2020 a dezembro de 2021 \\
\path{dummy_election_uncertainty} & outubro a dezembro dos anos eleitorais presidenciais 2002, 2006, 2010, 2014, 2018 e 2022 \\
\bottomrule
\end{longtable}

\subsubsection{Operações extraordinárias}

O script inspecionou \path{data/clean/auction_operations_clean.csv}. Os valores reais da
coluna \path{operation} encontrados foram:

\begin{itemize}[leftmargin=1.2cm]
    \item \path{VENDA};
    \item \path{COMPRA};
    \item \path{TROCA};
    \item \path{EXTRA COMPRA};
    \item \path{EXTRA VENDA}.
\end{itemize}

A partir disso, o script criou dummies mensais e medidas de intensidade:

\begin{itemize}[leftmargin=1.2cm]
    \item \path{dummy_extra_sale};
    \item \path{dummy_extra_purchase};
    \item \path{dummy_exchange_operation};
    \item \path{dummy_buyback_operation};
    \item \path{extra_sale_value_brl} e \path{extra_sale_value_pib};
    \item \path{extra_purchase_value_brl} e \path{extra_purchase_value_pib};
    \item \path{exchange_value_brl} e \path{exchange_value_pib};
    \item \path{buyback_value_brl} e \path{buyback_value_pib}.
\end{itemize}

Também foram criadas versões defasadas \path{lag1}. A normalização por PIB usa
\path{pib_12m_lag1_brl}.

\subsubsection{PAF}

Foi criado o template manual:
\begin{center}
\path{data/raw/paf_manual_targets.csv}.
\end{center}

O template inclui colunas para ano, datas de validade, limites mínimos e máximos
do PAF por indexador, percentual vincendo em 12 meses, prazo médio, estoque,
estratégias qualitativas, documento-fonte e notas. Como o arquivo está vazio, as
variáveis de PAF foram criadas no painel, mas permanecem ausentes até que o
template seja preenchido.

As variáveis programadas incluem:

\begin{itemize}[leftmargin=1.2cm]
    \item alvos médios: \path{paf_N_target}, \path{paf_F_target},
    \path{paf_R_target};
    \item larguras de banda: \path{paf_N_band_width}, \path{paf_F_band_width},
    \path{paf_R_band_width};
    \item gaps: \path{paf_gap_N_lag1}, \path{paf_gap_F_lag1},
    \path{paf_gap_R_lag1};
    \item distâncias aos limites inferiores e superiores;
    \item dummies de estar abaixo ou acima da banda;
    \item dummies de estratégia qualitativa \path{increase}, \path{decrease} e
    \path{stable} para \(N\), \(F\) e \(R\).
\end{itemize}

\subsubsection{Reserva de liquidez}

Foi criado o template manual:
\begin{center}
\path{data/raw/treasury_liquidity_reserve_manual.csv}.
\end{center}

O template contém:

\begin{itemize}[leftmargin=1.2cm]
    \item \path{month};
    \item \path{liquidity_reserve_brl};
    \item \path{liquidity_reserve_months_of_debt_service};
    \item \path{source_document};
    \item \path{notes}.
\end{itemize}

Como o template ainda está vazio, as variáveis de reserva de liquidez foram
criadas, mas ficam ausentes. Quando preenchido, o script criará:

\begin{itemize}[leftmargin=1.2cm]
    \item \path{liquidity_reserve_brl_lag1};
    \item \path{liquidity_reserve_pib_lag1};
    \item \path{liquidity_reserve_over_stock_lag1};
    \item \path{liquidity_reserve_over_maturity_12m_lag1}.
\end{itemize}

O valor contemporâneo não é usado como controle da emissão de \(t\); a reserva de
liquidez entra defasada.

\subsection{Listas finais de controles}

O script salva as listas em
\path{data/diagnostics/structural_control_lists.json}.

\subsubsection{Lista base}

A lista \path{structural_controls_base} gerada foi:

\begin{itemize}[leftmargin=1.2cm]
    \item \path{g_N_pib_lag1};
    \item \path{g_F_pib_lag1};
    \item \path{g_R_pib_lag1};
    \item \path{share_N_lag1};
    \item \path{share_F_lag1};
    \item \path{share_R_lag1};
    \item \path{paf_gap_N_lag1};
    \item \path{paf_gap_F_lag1};
    \item \path{paf_gap_R_lag1};
    \item \path{liquidity_reserve_pib_lag1};
    \item \path{m_total_NFR_pib_window_3m}.
\end{itemize}

\subsubsection{Lista estendida}

A lista \path{structural_controls_extended} inclui a lista base e adiciona:

\begin{itemize}[leftmargin=1.2cm]
    \item defasagens adicionais de emissão;
    \item médias móveis de emissão;
    \item dummies de crise e eleição;
    \item dummies e intensidades de operações extraordinárias;
    \item dummies qualitativas do PAF;
    \item distâncias às bandas do PAF;
    \item reserva de liquidez sobre estoque e vencimentos;
    \item janelas de vencimentos de 6 e 12 meses;
    \item janelas de vencimentos de 3 meses por classe.
\end{itemize}

Na execução atual, nenhum candidato esperado ficou ausente como coluna. Algumas
colunas, porém, estão completamente vazias porque dependem de templates manuais
ainda não preenchidos.

\subsection{Testes e diagnósticos executados}

\subsubsection{Execução local do script}

O script foi executado localmente com:
\begin{center}
\path{python3 src/build_structural_controls.py}.
\end{center}

Houve um problema inicial relacionado ao cache de fontes do \path{matplotlib} no
macOS. O script foi ajustado para criar \path{data/.mplconfig} e forçar o
backend \path{Agg}. Após esse ajuste, a execução terminou sem erro.

O resultado principal foi:

\begin{itemize}[leftmargin=1.2cm]
    \item painel final salvo em
    \path{data/clean/ftpl_monthly_panel_structural_controls.csv};
    \item dimensões: 559 linhas por 190 colunas;
    \item período: dezembro de 1979 a junho de 2026.
\end{itemize}

\subsubsection{Inspeção de colunas reais}

Antes da programação da Etapa A, foram inspecionados os nomes reais das colunas dos
CSVs. Isso confirmou que o painel mais recente tinha as variáveis normalizadas por
PIB e que a base de leilões tinha a coluna \path{operation}. Essa inspeção evitou
dependência de nomes assumidos na discussão.

\subsubsection{Teste das operações extraordinárias}

Foi verificada a distribuição real da coluna \path{operation} em
\path{auction_operations_clean.csv}. As categorias observadas permitiram mapear:

\begin{itemize}[leftmargin=1.2cm]
    \item \path{EXTRA VENDA} para venda extraordinária;
    \item \path{EXTRA COMPRA} para compra extraordinária;
    \item \path{TROCA} para operação de troca;
    \item \path{COMPRA} para recompra ou buyback.
\end{itemize}

Com isso, o script agregou operações por mês e gerou dummies e valores mensais.

\subsubsection{Diagnósticos salvos}

Foram salvos em \path{data/diagnostics/}:

\begin{longtable}{p{0.42\textwidth}p{0.48\textwidth}}
\toprule
\textbf{Arquivo} & \textbf{Conteúdo} \\
\midrule
\endhead
\path{structural_controls_missingness.csv} &
Tabela de valores ausentes das variáveis estruturais. \\
\path{structural_controls_descriptive_stats.csv} &
Estatísticas descritivas das variáveis estruturais numéricas. \\
\path{structural_controls_main_correlation.csv} &
Matriz de correlação das variáveis estruturais principais. \\
\path{structural_controls_alerts.csv} &
Alertas de variáveis com menos de 50\% de observações não ausentes ou variáveis
constantes/quase constantes. \\
\path{structural_control_lists.json} &
Listas finais de controles base e estendida, além de candidatos ausentes. \\
\bottomrule
\end{longtable}

Foi salvo em \path{data/figures/}:

\begin{itemize}[leftmargin=1.2cm]
    \item \path{structural_controls_main_correlation.png}.
\end{itemize}

O gráfico de PAF versus participações efetivas e o gráfico de reserva de liquidez
não foram gerados nesta execução porque os templates manuais correspondentes ainda
estão vazios.

\subsubsection{Alertas relevantes}

Os diagnósticos apontaram:

\begin{itemize}[leftmargin=1.2cm]
    \item \path{share_N_lag1}, \path{share_F_lag1} e \path{share_R_lag1} têm menos
    de 50\% de observações não ausentes, refletindo a cobertura mais curta dos
    estoques;
    \item na execução inicial, variáveis de PAF e reserva de liquidez estavam
    ausentes porque os templates ainda não haviam sido preenchidos; posteriormente
    esses templates foram completados e reintegrados ao painel;
    \item dummies de eventos e operações aparecem como quase constantes por serem
    variáveis binárias com poucos episódios, o que é esperado.
\end{itemize}

\section{Preenchimento auditável dos templates manuais}\label{app:templates}

Após a implementação inicial da Etapa A, os templates de PAF e reserva de liquidez
foram preenchidos com fontes oficiais. A lógica do pipeline principal não foi
alterada. O preenchimento foi feito em uma etapa auxiliar, preservando os nomes
originais das colunas dos templates e criando backups antes de qualquer escrita.

\subsection{Backups e segurança}

Antes do preenchimento, foram criados os backups:

\begin{itemize}[leftmargin=1.2cm]
    \item \path{data/raw/paf_manual_targets_backup_20260630.csv};
    \item \path{data/raw/treasury_liquidity_reserve_manual_backup_20260630.csv}.
\end{itemize}

Os templates originais estavam vazios, contendo apenas cabeçalhos. Os nomes das
colunas foram preservados exatamente.

\subsection{Script auxiliar de preenchimento}

Foi criado o script \path{src/fill_manual_templates_official_sources.py}. Este script é
separado do pipeline principal e tem apenas a função de preencher os templates
manuais e gerar arquivos de auditoria. Ele não altera a especificação econométrica
nem a lógica de construção das variáveis estruturais.

O script usa:

\begin{itemize}[leftmargin=1.2cm]
    \item \path{openpyxl} para ler anexos oficiais em \path{.xlsx};
    \item \path{pandas} para organizar e salvar os CSVs;
    \item fontes baixadas em \path{data/raw/paf_sources/} e
    \path{data/raw/liquidity_sources/}.
\end{itemize}

\subsection{Fontes oficiais para o PAF}

Para o PAF, foram usadas duas fontes oficiais:

\begin{enumerate}[leftmargin=1.2cm]
    \item O anexo oficial do RMD de maio de 2026, arquivo
    \path{Anexo RMD_Maio_26.xlsx}, aba \path{5.7 - Acompanhamento dos Planos
    Anuais de Financiamento da DPF}. Essa fonte contém dados estruturados dos PAFs
    de 2003 a 2022.
    \item Os PDFs oficiais do Plano Anual de Financiamento para 2023, 2024, 2025 e
    2026, baixados das páginas oficiais do Tesouro Transparente.
\end{enumerate}

Arquivos brutos preservados:

\begin{itemize}[leftmargin=1.2cm]
    \item \path{data/raw/paf_sources/PAF_2023.pdf};
    \item \path{data/raw/paf_sources/PAF_2024.pdf};
    \item \path{data/raw/paf_sources/PAF_2025.pdf};
    \item \path{data/raw/paf_sources/PAF_2026.pdf};
    \item \path{data/raw/paf_sources/Apresentacao_PAF_26_RAD25.pdf};
    \item páginas HTML intermediárias com os links oficiais das publicações.
\end{itemize}

O template preenchido é:
\begin{center}
\path{data/raw/paf_manual_targets.csv}.
\end{center}

Ele contém 24 linhas, cobrindo os anos de 2003 a 2026. As bandas numéricas foram
preenchidas para:

\begin{itemize}[leftmargin=1.2cm]
    \item participação de prefixados;
    \item participação de títulos remunerados pela Selic ou taxa flutuante;
    \item participação de títulos indexados a índices de preços;
    \item participação cambial;
    \item percentual vincendo em 12 meses;
    \item prazo médio;
    \item estoque da DPF.
\end{itemize}

As participações foram convertidas de percentual para proporção. Valores de estoque
reportados em R\$ bilhões foram convertidos para reais. As diretrizes qualitativas
foram mantidas como \path{unclear}, por conservadorismo: as fontes usadas trazem
bandas quantitativas claras, mas não permitem uma codificação anual inequívoca de
\path{increase}, \path{decrease} ou \path{stable} para todos os indexadores.

\subsection{Auditoria do PAF}

Foi criado:
\begin{center}
\path{data/diagnostics/paf_manual_targets_audit.csv}.
\end{center}

O arquivo contém 336 registros de auditoria, com as colunas:
\path{year}, \path{variable}, \path{value}, \path{source_document},
\path{source_url}, \path{source_file}, \path{page_or_table},
\path{extraction_note} e \path{confidence}.

A confiança foi classificada como:

\begin{itemize}[leftmargin=1.2cm]
    \item \path{high} para valores extraídos da aba estruturada oficial do RMD;
    \item \path{medium} para valores extraídos dos quadros claros dos PDFs oficiais
    do PAF.
\end{itemize}

\subsection{Fonte oficial para a reserva de liquidez}

A reserva de liquidez foi extraída do anexo oficial do RMD de maio de 2026:

\begin{center}
\path{data/raw/liquidity_sources/Anexo RMD_Maio_26.xlsx}.
\end{center}

A variável está na aba \path{6.1 - Reserva de Liquidez}. A tabela contém:

\begin{itemize}[leftmargin=1.2cm]
    \item \path{Mês};
    \item \path{Reserva de Liquidez}, em R\$ bilhões;
    \item \path{Valor Prudencial (3 meses)}, em R\$ bilhões;
    \item \path{Índice de Liquidez}, em meses.
\end{itemize}

O template preenchido é:
\begin{center}
\path{data/raw/treasury_liquidity_reserve_manual.csv}.
\end{center}

Ele contém 137 observações mensais, de janeiro de 2015 a maio de 2026. A reserva
foi convertida de R\$ bilhões para reais. O índice de liquidez foi mantido em meses
e armazenado em \path{liquidity_reserve_months_of_debt_service}.

\subsection{Auditoria da reserva de liquidez}

Foi criado:
\begin{center}
\path{data/diagnostics/treasury_liquidity_reserve_audit.csv}.
\end{center}

O arquivo contém 274 registros de auditoria: dois registros por mês, um para
\path{liquidity_reserve_brl} e outro para
\path{liquidity_reserve_months_of_debt_service}. A fonte é uma tabela estruturada
oficial, por isso a confiança foi classificada como \path{high}.

\subsection{Reintegração no painel estrutural}

Depois de preencher os templates, o script \path{src/build_structural_controls.py} foi
executado novamente. O painel atualizado foi salvo em:

\begin{center}
\path{data/clean/ftpl_monthly_panel_structural_controls.csv}.
\end{center}

O painel passou a ter 559 linhas e 192 colunas. A cobertura das principais
variáveis integradas ficou:

\begin{longtable}{p{0.44\textwidth}r p{0.10\textwidth}p{0.10\textwidth}}
\toprule
\textbf{Variável} & \textbf{Obs. não ausentes} & \textbf{Início} & \textbf{Fim} \\
\midrule
\endhead
\path{paf_N_target} & 282 & 2003-01 & 2026-06 \\
\path{paf_F_target} & 282 & 2003-01 & 2026-06 \\
\path{paf_R_target} & 282 & 2003-01 & 2026-06 \\
\path{paf_gap_N_lag1} & 104 & 2017-10 & 2026-05 \\
\path{paf_gap_F_lag1} & 104 & 2017-10 & 2026-05 \\
\path{paf_gap_R_lag1} & 104 & 2017-10 & 2026-05 \\
\path{liquidity_reserve_brl_lag1} & 137 & 2015-02 & 2026-06 \\
\path{liquidity_reserve_pib_lag1} & 136 & 2015-02 & 2026-05 \\
\path{liquidity_reserve_over_stock_lag1} & 104 & 2017-10 & 2026-05 \\
\path{liquidity_reserve_over_maturity_12m_lag1} & 114 & 2016-12 & 2026-05 \\
\bottomrule
\end{longtable}

Os gaps de PAF têm menos observações do que os alvos porque dependem das
participações efetivas defasadas \path{share_N_lag1}, \path{share_F_lag1} e
\path{share_R_lag1}, cuja cobertura começa apenas quando a base de estoque está
disponível.

\subsection{Diagnósticos e gráficos adicionais}

Além dos diagnósticos gerais da Etapa A, foram criados:

\begin{itemize}[leftmargin=1.2cm]
    \item \path{data/diagnostics/manual_templates_integration_summary.csv};
    \item \path{data/diagnostics/manual_templates_coverage_summary.csv};
    \item \path{data/figures/paf_effective_shares_vs_bands_manual_filled.png};
    \item \path{data/figures/liquidity_reserve_brl_billion.png};
    \item \path{data/figures/liquidity_reserve_pib_lag1_pct.png};
    \item \path{data/figures/liquidity_reserve_over_maturity_12m_lag1.png}.
\end{itemize}

Após o preenchimento, o gráfico de participações efetivas versus bandas do PAF e os
gráficos da reserva de liquidez passaram a ser gerados.

\section{Entregas do coautor}\label{app:coautor}

As entregas do coautor integradas ao pipeline estão descritas onde são usadas: a
base de notícia fiscal Prisma/Focus (Subseção~\ref{sec:fiscal_dados}) e a ETTJ
ANBIMA diária (Subseção~\ref{sec:ettj}). Os originais ficam em \path{DadosJoao/},
fora do Git, e as cópias usadas, com SHA-256, em \path{data/raw/}. Auditoria da
segunda entrega:

A pasta \path{DadosJoao/data/} foi auditada em 22/11/2026. Os pontos relevantes são:

\begin{itemize}[leftmargin=1.2cm]
    \item \path{stock_by_group_bcb_sgs.csv} não deve ser usado. Os valores brutos do
    SGS são idênticos aos nossos, mas a NTN-F foi inferida das participações
    percentuais 4177 e 4178, o que gera NTN-F negativa em 57 meses de 2000 a 2004,
    erro no estoque \(N\) em todos os meses (até 11,6\%, ou R\$ 7,4 bi) e erro de até
    1,7 ponto percentual na participação de \(N\);
    \item \path{fatores_curva_juros.xlsx} contém apenas os parâmetros \(\beta_1\) a
    \(\beta_4\) de Svensson, mensais de fim de mês, de 2010-01 a 2025-12, sem
    \(\lambda_1\) e \(\lambda_2\). Sem eles não é possível reconstruir os vértices
    nem a inflação implícita, e o arquivo não atende aos requisitos da Etapa B. Foram
    pedidos ao coautor os yields diários da curva zero-cupom da ANBIMA (prefixada,
    IPCA e inflação implícita), do início de cada série, no mínimo desde
    dezembro de 2009, até a data mais recente. Eles chegaram em seguida, em
    \path{data_yields.xlsx} (subseção~\ref{sec:ettj});
    \item \path{CDS_Brasil.xlsx} traz CDS de 5 e 2 anos de 2008-01 a 2026-06, sem fonte
    declarada e com meses de 2008 repetidos; fica apenas para robustez e não foi
    integrado;
    \item as planilhas brutas do Prisma de 2015 a 2023 cobrem parte do pedido
    anterior; as expectativas de 2024 em diante só existem como exportação do painel
    Qlik, ainda não preservada;
    \item os demais arquivos repetem versões anteriores da nossa própria cadeia.
\end{itemize}

\section{Histórico de correções}\label{app:hist}

Registro das correções que mudaram números já reportados. Estimativas anteriores
às datas abaixo não devem ser usadas.

\subsection{Leitura das séries do SGS (22/11/2026)}

A função \path{fetch_bcb_sgs} do notebook removia todos os pontos dos valores
retornados pelo SGS antes da conversão numérica, mas a API JSON do Banco Central usa
o ponto como separador decimal (por exemplo, \path{"638091.1"}). Cada série ficava
multiplicada por \(10^{d}\), onde \(d\) é o número de casas decimais:

\begin{center}
\small
\begin{tabular}{p{0.35\textwidth}p{0.12\textwidth}p{0.43\textwidth}}
\toprule
Série & Fator & Consequência \\
\midrule
PIB mensal (SGS 4380) & \(\times 10\) & todas as variáveis \path{*_pib} ficavam 10
vezes menores \\
Meta Selic (SGS 432) & \(\times 100\) & 4,50\% aparecia como 450 \\
Câmbio BRL/USD (SGS 3698) & \(\times 10^4\) & nível incorreto \\
IBC-Br & \(\times 10^5\) & nível incorreto \\
\bottomrule
\end{tabular}
\end{center}

Taxas de variação, logaritmos e participações não foram afetados; níveis e
combinações entre séries, sim. A correção mantém a conversão do formato brasileiro
apenas para valores com vírgula (fallback em CSV). Na mesma revisão, a coluna
\path{ibc_br_sa} passou da série 24363, sem ajuste sazonal, para a 24364, com
ajuste (a média da variação mensal de março caiu de +8,3\% para a faixa de
\(-0{,}3\%\) a \(+0{,}5\%\)). A cadeia inteira foi reexecutada; as colunas em reais
de emissão, estoque e vencimento ficaram idênticas. As saídas antigas estão em
\path{backups/sgs_decimal_fix_20261122/} (fora do Git).

\subsection{Etapa B: datas e busca de arquivos (22/11/2026)}

\begin{itemize}[leftmargin=1.2cm]
    \item A leitura de datas aplicava a regra de dia primeiro a toda data em texto,
    o que transformava ``2006-01-02'' em 1\textsuperscript{o} de fevereiro e deixava
    vazias as datas com dia acima de 12. Agora as datas ISO são lidas primeiro e as
    demais com dia primeiro.
    \item A busca de arquivos locais varria todo o repositório e encontrava cópias
    antigas em \path{backups/} e \path{DadosJoao/}. Agora ela se restringe a
    \path{data/raw/anbima/} e \path{data/anbima/} e ignora
    \path{data/raw/anbima/original/}.
\end{itemize}

\subsection{Datas da série realizada do Prisma}

Na série de resultado primário realizado entregue pelo coautor, os meses de julho
de 2020 a dezembro de 2021 estavam datados nos dias 2 a 13, o que fazia o merge por
mês descartar 18 surpresas. A surpresa foi recalculada e passou de 111 para 129
meses.

\subsection{Lotes conjuntos nos leilões (23/11/2026)}

O arquivo de leilões do Tesouro tem uma linha por título; quando vários vencimentos
dividem um lote conjunto, o lote é repetido em cada linha. A primeira versão do
diagnóstico de leilões contava por título e concluía que só 28\% dos leilões de
NTN-B vendiam o lote inteiro, de modo que a emissão de NTN-B seria em geral
limitada pela demanda. Contando por lote, a fração é 63\% (50\% tratando os grupos
ambíguos como lotes separados); em \(F\), 65\% em vez de 35\%. A persistência do
estado após um leilão não restritivo também era artefato (0,10 e 0,08 em \(F\) e
\(R\) por título, contra 0,41 e 0,44 por lote). As conclusões da
Seção~\ref{sec:ident} e a construção das surpresas de oferta usam a contagem por
lote.

\subsection{Contagem de colunas}

Uma versão anterior deste documento registrava 399 colunas no painel de estoque
antes da ETTJ; eram 400. As dimensões passaram a ficar só na tabela da
Seção~\ref{sec:pipeline}.

\subsection{Ambiente local}

Na primeira execução da Etapa A, o cache de fontes do \path{matplotlib} falhou no
macOS. Os scripts criam \path{data/.mplconfig} e forçam o backend \path{Agg}.

\end{document}
